% Pass: Opus 5 (claude-opus-5), 2026-09-12 - sheets 1-12, first pass,
% unchecked against the photographs by a human.
% Pass: Opus 5 (claude-opus-5), 2026-09-12 - sheets 13-24, first pass,
% unchecked against the photographs by a human.
%
% Re: 3 Wash Sq: a stenographer's spiral pad bound at the top, 51 photographed
% sheets, of which the first sitting reads the first twelve. It is neither a
% memorandum book like Battle of Wash Sq. nor a diary like Garrulities: it is a
% working draft. Page 1 is a list of the Village associations, the Sailors'
% Snug Harbor trustees and the officers of the Washington Square Association.
% Everything after it is one statement about the eviction of the tenants of
% 3 Washington Square North, written out three times over - a first draft in a
% fast hand (pages 2-3), a second and fuller one (pages 4-7), and a fair copy
% addressed to Mr. Whitelaw Reid of the Herald Tribune and dated Feb. 10, 47,
% which begins on the lower side of page 8 and runs to the end of page 11, and
% whose sides she numbers 2 to 7 in the head margin from page 9 on. A covering
% note to Reid, and a passage cancelled with a cross, stand between the second
% draft and the fair copy on pages 7 and 8. The same sentences therefore recur
% three and four times, and never twice the same: names are added, spelled
% differently, or dropped. Nothing here is merged, and each side is given as it
% stands.
%
% Each photograph shows two written sides of the pad - the leaf turned back
% above the spiral and the leaf below it - and the two run on continuously, the
% lower side of one photograph continuing into the upper side of the next. The
% Whitney's « Page N » names the photograph, not the side, so \sheet{} takes
% that number and both sides follow it in order.
%
% \entry{} is used at the four places where a date heads a passage: the O.P.A.
% hearing on pages 6 (twice) and 11, and the letter to Whitelaw Reid on pages 7
% and 8. The hearing is « Tues. Feb. 19" » three times and « Mon. A.M. Feb. 18" »
% once, and both stand as written.
%
% The second sitting reads pages 12 to 23, and the pad changes character in it.
% The Whitelaw Reid statement is finished on the upper side of page 12; below
% it the pad becomes a working file - names, telephone numbers, and the O.P.A.
% tax assessments of four houses on the Square. From page 13 what follows is
% not one statement written out again but six separate letters, each to a named
% person: John Morron of no. 6 and « dear Anita » on pages 13 and 14, a petition
% letter to Friedlander of 740 Park Avenue on pages 15 and 16, and three drafts
% to Mayor LaGuardia on pages 17-18, 19-20 and 21-22. The LaGuardia drafts share
% whole sentences with the Reid statement of the first sitting and with each
% other, and again nothing is merged. Page 20 carries the list of former tenants
% the drafts keep promising, in two columns; page 23 is a guest list for a party
% on Feb. 25, with ticks, crosses and telephone numbers against the names.
%
% The Friedlander letter is the one place the pad does not run forwards. Page 16
% is marked I and page 15's lower side is marked II, and page 16 ends with her
% own direction - « This begins on this page marked I + continues on page in
% back " II » - so the letter reads 16 then 15. Both are transcribed where they
% are bound, and the note on each says so.
%
% Two registers appear here that the first twelve sheets did not have. Red
% crayon marks pages 17, 20 and 23 - an asterisk and an underline beside
% « Painters require space », « 2 pages back » under the index on page 20, and
% dashes, underlines and asterisks down both columns of the tenant list and the
% guest list. And the word « Fiora » stands alone in the blank half of page 21
% in a much darker and blunter medium than the pencil around it; the hand is not
% placed.
%
% \entry{} is used at five places in this sitting where a date heads a passage:
% the two letters of Feb. 13 on pages 13 and 14, the LaGuardia draft of
% « Feb. 20" Thurs. » on page 21, and the party of « Tues. P.M. Feb. 25", 47 »
% on page 23.
\documentclass[11pt,a4paper]{article}
% The preamble shared by every transcription of the Hopper artist's ledgers.
%
% It rests on one idea, taken whole from the Grothendieck workbench this
% method comes from: **uncertainty compounds, it does not resolve.** A
% transcription that smooths over a doubtful figure has destroyed the only
% thing it was made for — a reader must be able to tell, at every point, what
% was on the leaf and what was supplied.
%
% A ledger sharpens that. In a manuscript of mathematics an invented word is a
% wrong word. Here an invented figure is a *sale that did not happen*, at a
% price nobody paid, to a buyer who never bought — and it will be cited,
% because a table looks like data in a way that prose does not.
%
% The same commands are recognised by `scripts/render.mjs`, which produces the
% site's reading view. A macro added here must be added there too, or rendering
% fails loudly, which is the wanted behaviour: a silently wrong rendering is
% worse than none.

% \ProvidesFile, not \ProvidesPackage: the transcriptions pull this in with
% % The preamble shared by every transcription of the Hopper artist's ledgers.
%
% It rests on one idea, taken whole from the Grothendieck workbench this
% method comes from: **uncertainty compounds, it does not resolve.** A
% transcription that smooths over a doubtful figure has destroyed the only
% thing it was made for — a reader must be able to tell, at every point, what
% was on the leaf and what was supplied.
%
% A ledger sharpens that. In a manuscript of mathematics an invented word is a
% wrong word. Here an invented figure is a *sale that did not happen*, at a
% price nobody paid, to a buyer who never bought — and it will be cited,
% because a table looks like data in a way that prose does not.
%
% The same commands are recognised by `scripts/render.mjs`, which produces the
% site's reading view. A macro added here must be added there too, or rendering
% fails loudly, which is the wanted behaviour: a silently wrong rendering is
% worse than none.

% \ProvidesFile, not \ProvidesPackage: the transcriptions pull this in with
% % The preamble shared by every transcription of the Hopper artist's ledgers.
%
% It rests on one idea, taken whole from the Grothendieck workbench this
% method comes from: **uncertainty compounds, it does not resolve.** A
% transcription that smooths over a doubtful figure has destroyed the only
% thing it was made for — a reader must be able to tell, at every point, what
% was on the leaf and what was supplied.
%
% A ledger sharpens that. In a manuscript of mathematics an invented word is a
% wrong word. Here an invented figure is a *sale that did not happen*, at a
% price nobody paid, to a buyer who never bought — and it will be cited,
% because a table looks like data in a way that prose does not.
%
% The same commands are recognised by `scripts/render.mjs`, which produces the
% site's reading view. A macro added here must be added there too, or rendering
% fails loudly, which is the wanted behaviour: a silently wrong rendering is
% worse than none.

% \ProvidesFile, not \ProvidesPackage: the transcriptions pull this in with
% \input{../preamble/hopper} rather than \usepackage, and \ProvidesPackage in
% an \input-ed file makes LaTeX warn on every compile that it "requested
% package `'" and got `hopper'. A warning nobody can act on is a warning
% everybody learns to scroll past, which is how the real ones get missed.
\ProvidesFile{hopper.sty}[2026/09/05 Transcriptions of the Hopper artist's ledgers]

\usepackage{array}
% longtable, not tabularx: a leaf of thirty-one exhibition lines must break
% across pages, and tabularx cannot be wrapped in an environment of our own —
% it scans the source for a literal \end{tabularx} and dies on \end{ledgertable}
% with « File ended while scanning use of \TX@get@body ». longtable has no such
% scan, and the wrapping column type below does the job tabularx's X would.
\usepackage{longtable}

% A wrapping column, `Y{0.3}` being three tenths of the measure.
%
% Prose columns must be declared this way and not as `l`. A table of `l`
% columns is set to its natural width, which can be wider than the page, and
% TeX issues **no warning at all** because nothing ever asked the row to fit:
% the first compile of Book I produced a page whose right-hand column ran off
% the paper with a clean log. A fixed-width column wraps, and if it still
% cannot fit, TeX says so — which is what `npm run pdf` then refuses to build
% over.
%
% Leave about 0.03 of the measure per column for the inter-column space: five
% columns can share 0.85, six can share 0.80. Getting it wrong is not a
% guessing game — `npm run pdf` reports the overfull box and how many points
% too wide it is.
\newcolumntype{Y}[1]{>{\raggedright\arraybackslash}p{#1\linewidth}}
\usepackage[english]{babel}
\usepackage{fontspec}
\usepackage{geometry}
\geometry{a4paper,margin=25mm,marginparwidth=28mm}
\usepackage{marginnote}
\usepackage{xcolor}
\usepackage{tikz}
% ulem, not soul. `soul`'s \st and \ul are built on a character-by-character
% scan that XeLaTeX's font machinery breaks: the document fails with an
% undefined control sequence rather than degrading. `normalem` stops it
% hijacking \emph, which the transcriptions use for underlining of Jo Hopper's
% own.
\usepackage[normalem]{ulem}
% ulem sets each underlined word in a box TeX can neither hyphenate nor break
% after, so a paragraph dense in \uncertain{} readings can reach a state where
% no legal breakpoint exists. \emergencystretch lets TeX loosen it on a third
% pass instead of overrunning: an overfull line in a transcription hides
% content.
\setlength{\emergencystretch}{3em}
\usepackage{hyperref}
\hypersetup{colorlinks=true,linkcolor=black,urlcolor=black,pdfborder={0 0 0}}

% --- Batch metadata ---------------------------------------------------------
% Taken as they stand by the rendering script, which shows them at the head of
% the reading view. They fix what the file claims to cover: without them a
% transcription floats unmoored in an archive of five hundred sheets.

\newcommand{\ledger}[1]{\gdef\hopLedger{#1}}
\newcommand{\ledgertitle}[1]{\gdef\hopLedgerTitle{#1}}
% A notebook of Josephine Hopper's, in place of a ledger: one file to a
% notebook, filed under transcripts/notebooks/, the argument the site's slug
% (« battle-of-wash-sq »). Sets the ledger to « notebooks » for everything
% that reads it, so the title block and the watermark behave as they do for a
% batch; the rendering checks the sheets against the notebook's own listing.
\newcommand{\notebook}[1]{\gdef\hopLedger{notebooks}\gdef\hopNotebook{#1}}
\gdef\hopNotebook{}
% The Whitney's accession number for the physical volume — 96.208 … 96.213.
% This is what a citation of the object cites.
\newcommand{\objectnumber}[1]{\gdef\hopObject{#1}}
% The batch this file covers. Leaving it out drops the « (batch N) » from the
% title block rather than printing a number that would be false.
\newcommand{\batch}[1]{\gdef\hopBatch{#1}}
\newcommand{\sheets}[2]{\gdef\hopFirst{#1}\gdef\hopLast{#2}}

% The Whitney's date range for the volume, copied verbatim from its resource
% record — never inferred here, never narrowed to the years this batch
% happens to cover.
\newcommand{\dating}[1]{\gdef\hopDating{#1}}

% The legal notice, printed in the title block and repeated as a diagonal
% watermark on every page of the PDF.
%
% This is not boilerplate and it is not optional. The ledgers are © Heirs of
% Josephine N. Hopper, licensed by the Artists Rights Society; the Whitney
% records the object rights as transferred to the Museum. A transcription
% reproduces Josephine Hopper's text. Every file must therefore say what it is
% on its own face, not only in a README that does not travel with the PDF.
% A file without this line gets no watermark, so omitting it is visible in the
% output rather than silent.
\newcommand{\watermark}[1]{\gdef\hopWatermark{#1}}

\gdef\hopLedger{?}\gdef\hopLedgerTitle{}\gdef\hopObject{}\gdef\hopBatch{}%
\gdef\hopFirst{?}\gdef\hopLast{?}\gdef\hopDating{}\gdef\hopWatermark{}

% --- Keywords ---------------------------------------------------------------
%
% Closing a transcription: the vocabulary under which to search for what these
% leaves record — dealers, exhibiting societies, media, places.
% `npm run manifest` extracts them to tag the whole ledger. This line is the
% single source of the tags; there is no tags file, so no tag can describe
% leaves nobody has read.
%
% It is the one thing in the file that is not on the paper, and it earns its
% place: a reader looking for Keppel needs some way in, and a tag written by
% whoever read the sheets is the only kind that cannot describe unread ones.
\newcommand{\keywords}[1]{\par\smallskip\noindent
  {\footnotesize\textsc{Keywords} --- \textit{#1}}\par}

% --- The anchor -------------------------------------------------------------

% The sheet begins here. Two arguments, and both are needed.
%
% #1 is the Whitney's **resource ref** — 16853 — which is the sheet's whole
% address: it is what the image URL is built from and what turns the facsimile
% as the transcript is scrolled. #2 is the number **written on the paper**, or
% empty where nobody wrote one, which is the case for every cover, flyleaf,
% index and loose insertion — about one sheet in nine.
%
% Keeping both is the whole of the numbering problem in this archive. Three
% sequences overlap: the Hoppers' own leaf numbers, the Whitney's order of
% photographs, and the resource refs, which are in no order at all. Only the
% first is citable and only the third is addressable, so a transcription that
% recorded one of them would be either uncitable or unopenable.
%
% `scripts/render.mjs` checks #2 against what `scripts/catalogue.mjs` parsed
% out of the Whitney's descriptor and fails on a disagreement, because a
% mistyped leaf number silently moves an entry to another year.
\newcommand{\sheet}[2]{%
  \marginnote{\footnotesize\color{gray!70}\textsf{\ifx\relax#2\relax---\else#2\fi}}%
  \label{sheet:#1}\ignorespaces}

% --- The critical apparatus -------------------------------------------------

% Illegible: never guessed. On a ledger leaf this is most often a figure, and
% a guessed figure is a fabricated transaction.
\newcommand{\ill}{\textcolor{red!60!black}{[\,\ldots\,]}}

% A doubtful reading: the word or figure is given, and flagged as uncertain.
%
% Underlining belongs to this macro and to nothing else. Jo Hopper underlines
% her own headings --- « Etchings », « Etching Plates », « Early oil » --- and
% those are set with \emph{} instead, because a reader must be able to tell a
% doubtful reading from her emphasis at a glance and the two cannot both be
% underlines. Which of them keeps the underline is decided by consequence: an
% emphasis shown as italic loses nothing, and a doubtful figure shown as a
% certain one loses everything.
\newcommand{\uncertain}[1]{\uline{#1}}

% An editorial addition — an expanded abbreviation, an implied dollar sign.
\newcommand{\add}[1]{\textcolor{blue!50!black}{[#1]}}

% Struck out in the book. Jo Hopper struck a great deal: a price revised, a
% buyer who withdrew, an exhibition that fell through. What she crossed out is
% part of the record and often the more interesting half of it.
\newcommand{\struck}[1]{\sout{#1}}

% The ink it is written in, where the ink says something.
%
% Book IV is a state machine and the colour is its status field: charges in
% black, receipts in red, the sums in pencil, a rule that holds from leaf 7 to
% the end of the volume. A transcription that records the words and the
% figures and nothing about their colour throws that field away, and the
% accounts then have to guess a receipt from the word « rec'd » --- which the
% volume stops writing on leaf 157. So the ink is recorded, per cell, as what
% is on the paper and not as what it means: \ink{red}{rec'd by check} says
% the words are red, and it is `scripts/accounts.mjs` that says red is money
% received. Black is the default and is not marked. The inks are the three
% the volume uses besides black; a fourth would be added here and in
% `scripts/render.mjs` and `scripts/tei.mjs` in the same commit.
% The file is \input, not \usepackage'd, so @ is not a letter here without saying so.
\makeatletter
\@namedef{hop@ink@red}{red!70!black}
\@namedef{hop@ink@pencil}{gray!55!black}
\@namedef{hop@ink@blue}{blue!55!black}
\newcommand{\ink}[2]{%
  \ifcsname hop@ink@#1\endcsname
    \textcolor{\csname hop@ink@#1\endcsname}{#2}%
  \else
    \PackageError{hopper}{\string\ink{#1}: unknown ink}{The inks are red, pencil and blue.}%
  \fi}
\makeatother

% A rule drawn across the money column above this row's figure.
%
% The writer rules off a run of items and writes their sum under the rule,
% and the rule is the whole of what says « this figure is a sum ». It stands
% at the head of the row the sum is on: \ruledoff & & 4355 & 00. On paper it
% is a line under the last two columns of the four a Book IV leaf has; if
% another volume ever rules a different column, the macro grows an argument.
\newcommand{\ruledoff}{\cline{3-4}}

% A diary entry begins. #1 is the date as she wrote it — « Mar. 12 »,
% « Sunday » — and #2 the date the transcriber assigns it, as a year, a month
% and a day (1947-03-12), or as much of one as the leaf allows. The first is
% what is on the paper; the second is the one editorial claim a diary needs
% that a ledger does not, because a ledger's leaf is its own address and an
% entry's date is its subject. Both are kept, apart, so the claim can be
% argued with. Used only in the notebooks; a ledger has no entries.
\newcommand{\entry}[2]{\par\medskip\noindent
  {\bfseries #1}\ifx\relax#2\relax\else\hspace{0.6em}{\footnotesize\color{gray!60!black}\textsf{#2}}\fi
  \par\nopagebreak\smallskip}

% The transcriber's note. Ours, not theirs.
\newcommand{\note}[1]{\par\smallskip{\small\color{gray!60!black}\textit{[#1]}}\par\smallskip}

% A marginal note **of theirs**, distinct from the body — and distinct from
% \note{} for exactly the reason \note{} exists.
\newcommand{\marginal}[1]{\par\smallskip\noindent\hspace{1em}%
  {\small\color{gray!60!black}#1}\par\smallskip}

% --- What is particular to this archive -------------------------------------

% Whose hand wrote it.
%
% This macro has no counterpart in the Grothendieck preamble, and it is the
% single most important thing in this one. These books were written by two
% people and annotated by more. Edward Hopper drew the work and wrote its
% title and size; Josephine Nivison Hopper wrote everything else — the
% descriptions, the anecdotes about people he refused to discuss, the prices,
% the buyers, the columns; and later hands, curatorial and unidentified, added
% to the leaves after 1967.
%
% A transcription that flattens the three has destroyed the archive's chief
% documentary value, which is that it records *two* working lives and the
% division of labour between them. #1 is `edward`, `jo`, `later` or
% `unidentified`; `unidentified` is the right answer far more often than a
% transcriber's confidence suggests, and it is not a failure to use it.
\newcommand{\hand}[2]{%
  \ifx\relax#1\relax#2\else
    \textcolor{gray!55!black}{\footnotesize\textsf{[#1]}}\,#2%
  \fi}

% Edward Hopper's ink record sketch stands here.
%
% The argument is **what is inscribed on or beside the sketch** — a title, a
% dimension, a note — and never a description of the image. Describing the
% drawing would be writing a new document: the sketch is on the site, one pane
% away, at the resolution the Whitney publishes, and prose about it competes
% with looking at it. Leave the argument empty where nothing is written.
% The mark is drawn rather than set from a symbol font: there is no
% mathematics anywhere in this archive, so no maths package is loaded, and
% $\square$ would be an undefined control sequence — which is exactly how it
% failed the first time the PDFs were compiled.
\newcommand{\sketchmark}{\raisebox{0.15ex}{\color{gray!50!black}%
  \tikz[baseline]{\draw[line width=0.4pt] (0,0) rectangle (1.1ex,1.1ex);}}}
\newcommand{\sketch}[1]{\par\smallskip\noindent
  {\small\color{gray!50!black}\sketchmark~\textsf{ink sketch}%
   \ifx\relax#1\relax\else\ --- #1\fi}\par\smallskip}

% Something printed and pasted to the leaf — a reproduction cut from a
% catalogue, a review, a dealer's label. The argument transcribes its printed
% text. These are the only places in the archive where the words are nobody's
% handwriting, and they date the leaf, which handwriting does not.
\newcommand{\clipping}[1]{\par\smallskip\noindent
  {\small\color{gray!50!black}\textsf{[clipping]}\ #1}\par\smallskip}

% A work's block opens here, under the title **exactly as the leaf gives it** —
% « Les Deux Pigeon » and « Les Poillus » stay misspelled, because the
% misspelling is Hopper's and is how the entry is found.
\newcommand{\work}[1]{\par\medskip\noindent{\bfseries #1}\par\nopagebreak\smallskip}

% --- The ruled table --------------------------------------------------------
%
% Jo Hopper ruled these columns herself and kept them for forty years; the
% table is not a presentation of the content, it *is* the content. So it is
% transcribed as a table and rendered as one, on screen and on paper, from one
% source.
%
% #1 is the column specification; #2 is the header row, `&`-separated like any
% other. A cell she left blank is left blank; a cell that cannot be read takes
% \ill{}.
%
% **Use `Y{...}` for any column that can hold a sentence** — an exhibition, a
% dealer's terms, a note — and `l`, `r` or `c` only for dates, figures and
% short names. See the note on \newcolumntype{Y} above for why that is not a
% matter of taste.
%
% The header row is emitted **bare**, between rules, and is deliberately not
% wrapped in \textbf{}. It cannot be: an `&` inside a braced group is not an
% alignment tab, so `\textbf{Date & Exhibitions}` fails at the first header
% with « Forbidden control sequence found while scanning use of
% \check@nocorr@ » — which is TeX saying, at some distance, that the tabs are
% in the wrong place. Rules carry the header instead, which is how a scholarly
% table sets one anyway. (The reading view has no such constraint and renders
% the same row as <th>.)
\newenvironment{ledgertable}[2]
  {\par\smallskip\small
   \begin{longtable}{#1}
   \hline
   #2 \\
   \hline
   \endhead}
  {\end{longtable}\par\smallskip}

% --- The appendix of notes --------------------------------------------------
%
% Everything above the appendix is on the paper. Nothing in it is.
%
% A note says what a named source says about a work these leaves record --- what
% the holding museum measures the canvas at, which other leaf carries the other
% half of a transaction --- and it is the only prose in the file that a reader
% cannot check against the photograph. On the site that is handled by putting
% the note beside a link to its source; on paper a link is not enough, because
% a PDF outlives the site it was downloaded from and travels away from it. So
% every claim here prints its source in full, name and address both, and every
% note prints the day it was written and what wrote it.
%
% The appendix is assembled by `scripts/pdf.mjs` from
% `src/content/work-notes.json` at compile time and is in no `.tex` file. That
% is the point of it. The transcription is the edition; a note is a later and
% weaker kind of writing *about* the edition, and keeping the two in one source
% file would eventually put a sentence somebody inferred inside the document
% whose whole claim on a reader is that nothing in it was inferred.
%
% Only the notes on works this batch's leaves name are printed, so the appendix
% follows the batch rather than repeating the archive, and a batch that has
% earned no notes yet gets no appendix at all --- which is a true statement
% about how much has been checked, and the same statement `work-notes.json`
% makes by being nearly empty.

% A URL is the one string in this archive that can run a hundred characters
% without a space in it, and an overfull box fails the build here because in a
% transcription an overfull box hides content. So the appendix tells TeX where
% a URL may break and sets its apparatus ragged right, rather than leaving the
% line to overrun in a file nobody recompiles.
\def\UrlBreaks{\do\/\do\-\do\.\do\_\do\?\do\&\do\=\do\+}
\Urlmuskip=0mu plus 1mu

\newcommand{\notesappendix}{\clearpage
  \par\noindent{\large\bfseries Notes on the works}\par\smallskip
  \noindent{\footnotesize\color{gray!60!black}\raggedright
    Not on the paper, and not part of the transcription. Each note below
    concerns a work named on a leaf of this batch and says what a named source
    says about it --- a museum's own catalogue record, or another leaf of these
    ledgers. Every sentence carries the source that states it, printed with its
    address so that it can be followed from a copy of this file alone. Where a
    ledger and a museum disagree the disagreement is printed rather than
    resolved.\par}\par\medskip}

% One work: the title as the leaves give it, then the note.
\newcommand{\worknote}[2]{\par\medskip\noindent{\bfseries #1}\par\nopagebreak
  \smallskip\noindent #2\par}

% One claim under it: what is asserted, and who asserts it.
\newcommand{\workclaim}[2]{\par\smallskip\noindent
  {\footnotesize\color{gray!60!black}\raggedright\hangindent=1.6em\hangafter=0
   \hspace{1em}#1\ --- #2\par}}

% Who wrote the note and when, closing the block. A note is a model's reading
% of retrieved records and is weaker evidence than anything above it; the file
% records that rather than letting the appendix pass for the edition.
\newcommand{\worknoteby}[2]{\par\smallskip\noindent
  {\scriptsize\color{gray!55!black}\hspace{1em}Written #1 by #2.\par}\par\smallskip}

% --- Title ------------------------------------------------------------------

% The watermark is drawn on the shipout background so it sits under the text of
% every page, diagonal and light enough to leave the record readable.
\AddToHook{shipout/background}{%
  \ifx\hopWatermark\empty\else
    \begin{tikzpicture}[remember picture, overlay]
      \node[rotate=52, scale=3.4, align=center, text=gray!18,
            font=\sffamily\bfseries]
        at (current page.center) {\hopWatermark};
    \end{tikzpicture}%
  \fi}

\AtBeginDocument{%
  \begin{center}
    {\large\bfseries Edward and Josephine Hopper --- \hopLedgerTitle}\\[2pt]
    \ifx\hopObject\empty\else
      {\small Whitney Museum of American Art, \hopObject}\\[4pt]
    \fi
    {\small sheets \hopFirst--\hopLast
      \ifx\hopBatch\empty\else\ (batch \hopBatch)\fi}\\[2pt]
    \ifx\hopDating\empty\else
      {\small\color{gray!60!black}The Whitney's dating: \hopDating}\\[2pt]
    \fi
    \ifx\hopWatermark\empty\else
      {\small\bfseries\color{red!45!gray}\hopWatermark}\\[2pt]
    \fi
    {\footnotesize\color{gray!60!black}Transcribed from the digitised sheets published by the
     Whitney Museum of American Art.\\
     \textcopyright{} Heirs of Josephine N. Hopper, licensed by Artists Rights Society (ARS),
     New York.\\
     Uncertain readings are underlined, illegible passages marked [\,\ldots\,],
     editorial additions bracketed.}
  \end{center}
  \vspace{1em}}

\endinput
 rather than \usepackage, and \ProvidesPackage in
% an \input-ed file makes LaTeX warn on every compile that it "requested
% package `'" and got `hopper'. A warning nobody can act on is a warning
% everybody learns to scroll past, which is how the real ones get missed.
\ProvidesFile{hopper.sty}[2026/09/05 Transcriptions of the Hopper artist's ledgers]

\usepackage{array}
% longtable, not tabularx: a leaf of thirty-one exhibition lines must break
% across pages, and tabularx cannot be wrapped in an environment of our own —
% it scans the source for a literal \end{tabularx} and dies on \end{ledgertable}
% with « File ended while scanning use of \TX@get@body ». longtable has no such
% scan, and the wrapping column type below does the job tabularx's X would.
\usepackage{longtable}

% A wrapping column, `Y{0.3}` being three tenths of the measure.
%
% Prose columns must be declared this way and not as `l`. A table of `l`
% columns is set to its natural width, which can be wider than the page, and
% TeX issues **no warning at all** because nothing ever asked the row to fit:
% the first compile of Book I produced a page whose right-hand column ran off
% the paper with a clean log. A fixed-width column wraps, and if it still
% cannot fit, TeX says so — which is what `npm run pdf` then refuses to build
% over.
%
% Leave about 0.03 of the measure per column for the inter-column space: five
% columns can share 0.85, six can share 0.80. Getting it wrong is not a
% guessing game — `npm run pdf` reports the overfull box and how many points
% too wide it is.
\newcolumntype{Y}[1]{>{\raggedright\arraybackslash}p{#1\linewidth}}
\usepackage[english]{babel}
\usepackage{fontspec}
\usepackage{geometry}
\geometry{a4paper,margin=25mm,marginparwidth=28mm}
\usepackage{marginnote}
\usepackage{xcolor}
\usepackage{tikz}
% ulem, not soul. `soul`'s \st and \ul are built on a character-by-character
% scan that XeLaTeX's font machinery breaks: the document fails with an
% undefined control sequence rather than degrading. `normalem` stops it
% hijacking \emph, which the transcriptions use for underlining of Jo Hopper's
% own.
\usepackage[normalem]{ulem}
% ulem sets each underlined word in a box TeX can neither hyphenate nor break
% after, so a paragraph dense in \uncertain{} readings can reach a state where
% no legal breakpoint exists. \emergencystretch lets TeX loosen it on a third
% pass instead of overrunning: an overfull line in a transcription hides
% content.
\setlength{\emergencystretch}{3em}
\usepackage{hyperref}
\hypersetup{colorlinks=true,linkcolor=black,urlcolor=black,pdfborder={0 0 0}}

% --- Batch metadata ---------------------------------------------------------
% Taken as they stand by the rendering script, which shows them at the head of
% the reading view. They fix what the file claims to cover: without them a
% transcription floats unmoored in an archive of five hundred sheets.

\newcommand{\ledger}[1]{\gdef\hopLedger{#1}}
\newcommand{\ledgertitle}[1]{\gdef\hopLedgerTitle{#1}}
% A notebook of Josephine Hopper's, in place of a ledger: one file to a
% notebook, filed under transcripts/notebooks/, the argument the site's slug
% (« battle-of-wash-sq »). Sets the ledger to « notebooks » for everything
% that reads it, so the title block and the watermark behave as they do for a
% batch; the rendering checks the sheets against the notebook's own listing.
\newcommand{\notebook}[1]{\gdef\hopLedger{notebooks}\gdef\hopNotebook{#1}}
\gdef\hopNotebook{}
% The Whitney's accession number for the physical volume — 96.208 … 96.213.
% This is what a citation of the object cites.
\newcommand{\objectnumber}[1]{\gdef\hopObject{#1}}
% The batch this file covers. Leaving it out drops the « (batch N) » from the
% title block rather than printing a number that would be false.
\newcommand{\batch}[1]{\gdef\hopBatch{#1}}
\newcommand{\sheets}[2]{\gdef\hopFirst{#1}\gdef\hopLast{#2}}

% The Whitney's date range for the volume, copied verbatim from its resource
% record — never inferred here, never narrowed to the years this batch
% happens to cover.
\newcommand{\dating}[1]{\gdef\hopDating{#1}}

% The legal notice, printed in the title block and repeated as a diagonal
% watermark on every page of the PDF.
%
% This is not boilerplate and it is not optional. The ledgers are © Heirs of
% Josephine N. Hopper, licensed by the Artists Rights Society; the Whitney
% records the object rights as transferred to the Museum. A transcription
% reproduces Josephine Hopper's text. Every file must therefore say what it is
% on its own face, not only in a README that does not travel with the PDF.
% A file without this line gets no watermark, so omitting it is visible in the
% output rather than silent.
\newcommand{\watermark}[1]{\gdef\hopWatermark{#1}}

\gdef\hopLedger{?}\gdef\hopLedgerTitle{}\gdef\hopObject{}\gdef\hopBatch{}%
\gdef\hopFirst{?}\gdef\hopLast{?}\gdef\hopDating{}\gdef\hopWatermark{}

% --- Keywords ---------------------------------------------------------------
%
% Closing a transcription: the vocabulary under which to search for what these
% leaves record — dealers, exhibiting societies, media, places.
% `npm run manifest` extracts them to tag the whole ledger. This line is the
% single source of the tags; there is no tags file, so no tag can describe
% leaves nobody has read.
%
% It is the one thing in the file that is not on the paper, and it earns its
% place: a reader looking for Keppel needs some way in, and a tag written by
% whoever read the sheets is the only kind that cannot describe unread ones.
\newcommand{\keywords}[1]{\par\smallskip\noindent
  {\footnotesize\textsc{Keywords} --- \textit{#1}}\par}

% --- The anchor -------------------------------------------------------------

% The sheet begins here. Two arguments, and both are needed.
%
% #1 is the Whitney's **resource ref** — 16853 — which is the sheet's whole
% address: it is what the image URL is built from and what turns the facsimile
% as the transcript is scrolled. #2 is the number **written on the paper**, or
% empty where nobody wrote one, which is the case for every cover, flyleaf,
% index and loose insertion — about one sheet in nine.
%
% Keeping both is the whole of the numbering problem in this archive. Three
% sequences overlap: the Hoppers' own leaf numbers, the Whitney's order of
% photographs, and the resource refs, which are in no order at all. Only the
% first is citable and only the third is addressable, so a transcription that
% recorded one of them would be either uncitable or unopenable.
%
% `scripts/render.mjs` checks #2 against what `scripts/catalogue.mjs` parsed
% out of the Whitney's descriptor and fails on a disagreement, because a
% mistyped leaf number silently moves an entry to another year.
\newcommand{\sheet}[2]{%
  \marginnote{\footnotesize\color{gray!70}\textsf{\ifx\relax#2\relax---\else#2\fi}}%
  \label{sheet:#1}\ignorespaces}

% --- The critical apparatus -------------------------------------------------

% Illegible: never guessed. On a ledger leaf this is most often a figure, and
% a guessed figure is a fabricated transaction.
\newcommand{\ill}{\textcolor{red!60!black}{[\,\ldots\,]}}

% A doubtful reading: the word or figure is given, and flagged as uncertain.
%
% Underlining belongs to this macro and to nothing else. Jo Hopper underlines
% her own headings --- « Etchings », « Etching Plates », « Early oil » --- and
% those are set with \emph{} instead, because a reader must be able to tell a
% doubtful reading from her emphasis at a glance and the two cannot both be
% underlines. Which of them keeps the underline is decided by consequence: an
% emphasis shown as italic loses nothing, and a doubtful figure shown as a
% certain one loses everything.
\newcommand{\uncertain}[1]{\uline{#1}}

% An editorial addition — an expanded abbreviation, an implied dollar sign.
\newcommand{\add}[1]{\textcolor{blue!50!black}{[#1]}}

% Struck out in the book. Jo Hopper struck a great deal: a price revised, a
% buyer who withdrew, an exhibition that fell through. What she crossed out is
% part of the record and often the more interesting half of it.
\newcommand{\struck}[1]{\sout{#1}}

% The ink it is written in, where the ink says something.
%
% Book IV is a state machine and the colour is its status field: charges in
% black, receipts in red, the sums in pencil, a rule that holds from leaf 7 to
% the end of the volume. A transcription that records the words and the
% figures and nothing about their colour throws that field away, and the
% accounts then have to guess a receipt from the word « rec'd » --- which the
% volume stops writing on leaf 157. So the ink is recorded, per cell, as what
% is on the paper and not as what it means: \ink{red}{rec'd by check} says
% the words are red, and it is `scripts/accounts.mjs` that says red is money
% received. Black is the default and is not marked. The inks are the three
% the volume uses besides black; a fourth would be added here and in
% `scripts/render.mjs` and `scripts/tei.mjs` in the same commit.
% The file is \input, not \usepackage'd, so @ is not a letter here without saying so.
\makeatletter
\@namedef{hop@ink@red}{red!70!black}
\@namedef{hop@ink@pencil}{gray!55!black}
\@namedef{hop@ink@blue}{blue!55!black}
\newcommand{\ink}[2]{%
  \ifcsname hop@ink@#1\endcsname
    \textcolor{\csname hop@ink@#1\endcsname}{#2}%
  \else
    \PackageError{hopper}{\string\ink{#1}: unknown ink}{The inks are red, pencil and blue.}%
  \fi}
\makeatother

% A rule drawn across the money column above this row's figure.
%
% The writer rules off a run of items and writes their sum under the rule,
% and the rule is the whole of what says « this figure is a sum ». It stands
% at the head of the row the sum is on: \ruledoff & & 4355 & 00. On paper it
% is a line under the last two columns of the four a Book IV leaf has; if
% another volume ever rules a different column, the macro grows an argument.
\newcommand{\ruledoff}{\cline{3-4}}

% A diary entry begins. #1 is the date as she wrote it — « Mar. 12 »,
% « Sunday » — and #2 the date the transcriber assigns it, as a year, a month
% and a day (1947-03-12), or as much of one as the leaf allows. The first is
% what is on the paper; the second is the one editorial claim a diary needs
% that a ledger does not, because a ledger's leaf is its own address and an
% entry's date is its subject. Both are kept, apart, so the claim can be
% argued with. Used only in the notebooks; a ledger has no entries.
\newcommand{\entry}[2]{\par\medskip\noindent
  {\bfseries #1}\ifx\relax#2\relax\else\hspace{0.6em}{\footnotesize\color{gray!60!black}\textsf{#2}}\fi
  \par\nopagebreak\smallskip}

% The transcriber's note. Ours, not theirs.
\newcommand{\note}[1]{\par\smallskip{\small\color{gray!60!black}\textit{[#1]}}\par\smallskip}

% A marginal note **of theirs**, distinct from the body — and distinct from
% \note{} for exactly the reason \note{} exists.
\newcommand{\marginal}[1]{\par\smallskip\noindent\hspace{1em}%
  {\small\color{gray!60!black}#1}\par\smallskip}

% --- What is particular to this archive -------------------------------------

% Whose hand wrote it.
%
% This macro has no counterpart in the Grothendieck preamble, and it is the
% single most important thing in this one. These books were written by two
% people and annotated by more. Edward Hopper drew the work and wrote its
% title and size; Josephine Nivison Hopper wrote everything else — the
% descriptions, the anecdotes about people he refused to discuss, the prices,
% the buyers, the columns; and later hands, curatorial and unidentified, added
% to the leaves after 1967.
%
% A transcription that flattens the three has destroyed the archive's chief
% documentary value, which is that it records *two* working lives and the
% division of labour between them. #1 is `edward`, `jo`, `later` or
% `unidentified`; `unidentified` is the right answer far more often than a
% transcriber's confidence suggests, and it is not a failure to use it.
\newcommand{\hand}[2]{%
  \ifx\relax#1\relax#2\else
    \textcolor{gray!55!black}{\footnotesize\textsf{[#1]}}\,#2%
  \fi}

% Edward Hopper's ink record sketch stands here.
%
% The argument is **what is inscribed on or beside the sketch** — a title, a
% dimension, a note — and never a description of the image. Describing the
% drawing would be writing a new document: the sketch is on the site, one pane
% away, at the resolution the Whitney publishes, and prose about it competes
% with looking at it. Leave the argument empty where nothing is written.
% The mark is drawn rather than set from a symbol font: there is no
% mathematics anywhere in this archive, so no maths package is loaded, and
% $\square$ would be an undefined control sequence — which is exactly how it
% failed the first time the PDFs were compiled.
\newcommand{\sketchmark}{\raisebox{0.15ex}{\color{gray!50!black}%
  \tikz[baseline]{\draw[line width=0.4pt] (0,0) rectangle (1.1ex,1.1ex);}}}
\newcommand{\sketch}[1]{\par\smallskip\noindent
  {\small\color{gray!50!black}\sketchmark~\textsf{ink sketch}%
   \ifx\relax#1\relax\else\ --- #1\fi}\par\smallskip}

% Something printed and pasted to the leaf — a reproduction cut from a
% catalogue, a review, a dealer's label. The argument transcribes its printed
% text. These are the only places in the archive where the words are nobody's
% handwriting, and they date the leaf, which handwriting does not.
\newcommand{\clipping}[1]{\par\smallskip\noindent
  {\small\color{gray!50!black}\textsf{[clipping]}\ #1}\par\smallskip}

% A work's block opens here, under the title **exactly as the leaf gives it** —
% « Les Deux Pigeon » and « Les Poillus » stay misspelled, because the
% misspelling is Hopper's and is how the entry is found.
\newcommand{\work}[1]{\par\medskip\noindent{\bfseries #1}\par\nopagebreak\smallskip}

% --- The ruled table --------------------------------------------------------
%
% Jo Hopper ruled these columns herself and kept them for forty years; the
% table is not a presentation of the content, it *is* the content. So it is
% transcribed as a table and rendered as one, on screen and on paper, from one
% source.
%
% #1 is the column specification; #2 is the header row, `&`-separated like any
% other. A cell she left blank is left blank; a cell that cannot be read takes
% \ill{}.
%
% **Use `Y{...}` for any column that can hold a sentence** — an exhibition, a
% dealer's terms, a note — and `l`, `r` or `c` only for dates, figures and
% short names. See the note on \newcolumntype{Y} above for why that is not a
% matter of taste.
%
% The header row is emitted **bare**, between rules, and is deliberately not
% wrapped in \textbf{}. It cannot be: an `&` inside a braced group is not an
% alignment tab, so `\textbf{Date & Exhibitions}` fails at the first header
% with « Forbidden control sequence found while scanning use of
% \check@nocorr@ » — which is TeX saying, at some distance, that the tabs are
% in the wrong place. Rules carry the header instead, which is how a scholarly
% table sets one anyway. (The reading view has no such constraint and renders
% the same row as <th>.)
\newenvironment{ledgertable}[2]
  {\par\smallskip\small
   \begin{longtable}{#1}
   \hline
   #2 \\
   \hline
   \endhead}
  {\end{longtable}\par\smallskip}

% --- The appendix of notes --------------------------------------------------
%
% Everything above the appendix is on the paper. Nothing in it is.
%
% A note says what a named source says about a work these leaves record --- what
% the holding museum measures the canvas at, which other leaf carries the other
% half of a transaction --- and it is the only prose in the file that a reader
% cannot check against the photograph. On the site that is handled by putting
% the note beside a link to its source; on paper a link is not enough, because
% a PDF outlives the site it was downloaded from and travels away from it. So
% every claim here prints its source in full, name and address both, and every
% note prints the day it was written and what wrote it.
%
% The appendix is assembled by `scripts/pdf.mjs` from
% `src/content/work-notes.json` at compile time and is in no `.tex` file. That
% is the point of it. The transcription is the edition; a note is a later and
% weaker kind of writing *about* the edition, and keeping the two in one source
% file would eventually put a sentence somebody inferred inside the document
% whose whole claim on a reader is that nothing in it was inferred.
%
% Only the notes on works this batch's leaves name are printed, so the appendix
% follows the batch rather than repeating the archive, and a batch that has
% earned no notes yet gets no appendix at all --- which is a true statement
% about how much has been checked, and the same statement `work-notes.json`
% makes by being nearly empty.

% A URL is the one string in this archive that can run a hundred characters
% without a space in it, and an overfull box fails the build here because in a
% transcription an overfull box hides content. So the appendix tells TeX where
% a URL may break and sets its apparatus ragged right, rather than leaving the
% line to overrun in a file nobody recompiles.
\def\UrlBreaks{\do\/\do\-\do\.\do\_\do\?\do\&\do\=\do\+}
\Urlmuskip=0mu plus 1mu

\newcommand{\notesappendix}{\clearpage
  \par\noindent{\large\bfseries Notes on the works}\par\smallskip
  \noindent{\footnotesize\color{gray!60!black}\raggedright
    Not on the paper, and not part of the transcription. Each note below
    concerns a work named on a leaf of this batch and says what a named source
    says about it --- a museum's own catalogue record, or another leaf of these
    ledgers. Every sentence carries the source that states it, printed with its
    address so that it can be followed from a copy of this file alone. Where a
    ledger and a museum disagree the disagreement is printed rather than
    resolved.\par}\par\medskip}

% One work: the title as the leaves give it, then the note.
\newcommand{\worknote}[2]{\par\medskip\noindent{\bfseries #1}\par\nopagebreak
  \smallskip\noindent #2\par}

% One claim under it: what is asserted, and who asserts it.
\newcommand{\workclaim}[2]{\par\smallskip\noindent
  {\footnotesize\color{gray!60!black}\raggedright\hangindent=1.6em\hangafter=0
   \hspace{1em}#1\ --- #2\par}}

% Who wrote the note and when, closing the block. A note is a model's reading
% of retrieved records and is weaker evidence than anything above it; the file
% records that rather than letting the appendix pass for the edition.
\newcommand{\worknoteby}[2]{\par\smallskip\noindent
  {\scriptsize\color{gray!55!black}\hspace{1em}Written #1 by #2.\par}\par\smallskip}

% --- Title ------------------------------------------------------------------

% The watermark is drawn on the shipout background so it sits under the text of
% every page, diagonal and light enough to leave the record readable.
\AddToHook{shipout/background}{%
  \ifx\hopWatermark\empty\else
    \begin{tikzpicture}[remember picture, overlay]
      \node[rotate=52, scale=3.4, align=center, text=gray!18,
            font=\sffamily\bfseries]
        at (current page.center) {\hopWatermark};
    \end{tikzpicture}%
  \fi}

\AtBeginDocument{%
  \begin{center}
    {\large\bfseries Edward and Josephine Hopper --- \hopLedgerTitle}\\[2pt]
    \ifx\hopObject\empty\else
      {\small Whitney Museum of American Art, \hopObject}\\[4pt]
    \fi
    {\small sheets \hopFirst--\hopLast
      \ifx\hopBatch\empty\else\ (batch \hopBatch)\fi}\\[2pt]
    \ifx\hopDating\empty\else
      {\small\color{gray!60!black}The Whitney's dating: \hopDating}\\[2pt]
    \fi
    \ifx\hopWatermark\empty\else
      {\small\bfseries\color{red!45!gray}\hopWatermark}\\[2pt]
    \fi
    {\footnotesize\color{gray!60!black}Transcribed from the digitised sheets published by the
     Whitney Museum of American Art.\\
     \textcopyright{} Heirs of Josephine N. Hopper, licensed by Artists Rights Society (ARS),
     New York.\\
     Uncertain readings are underlined, illegible passages marked [\,\ldots\,],
     editorial additions bracketed.}
  \end{center}
  \vspace{1em}}

\endinput
 rather than \usepackage, and \ProvidesPackage in
% an \input-ed file makes LaTeX warn on every compile that it "requested
% package `'" and got `hopper'. A warning nobody can act on is a warning
% everybody learns to scroll past, which is how the real ones get missed.
\ProvidesFile{hopper.sty}[2026/09/05 Transcriptions of the Hopper artist's ledgers]

\usepackage{array}
% longtable, not tabularx: a leaf of thirty-one exhibition lines must break
% across pages, and tabularx cannot be wrapped in an environment of our own —
% it scans the source for a literal \end{tabularx} and dies on \end{ledgertable}
% with « File ended while scanning use of \TX@get@body ». longtable has no such
% scan, and the wrapping column type below does the job tabularx's X would.
\usepackage{longtable}

% A wrapping column, `Y{0.3}` being three tenths of the measure.
%
% Prose columns must be declared this way and not as `l`. A table of `l`
% columns is set to its natural width, which can be wider than the page, and
% TeX issues **no warning at all** because nothing ever asked the row to fit:
% the first compile of Book I produced a page whose right-hand column ran off
% the paper with a clean log. A fixed-width column wraps, and if it still
% cannot fit, TeX says so — which is what `npm run pdf` then refuses to build
% over.
%
% Leave about 0.03 of the measure per column for the inter-column space: five
% columns can share 0.85, six can share 0.80. Getting it wrong is not a
% guessing game — `npm run pdf` reports the overfull box and how many points
% too wide it is.
\newcolumntype{Y}[1]{>{\raggedright\arraybackslash}p{#1\linewidth}}
\usepackage[english]{babel}
\usepackage{fontspec}
\usepackage{geometry}
\geometry{a4paper,margin=25mm,marginparwidth=28mm}
\usepackage{marginnote}
\usepackage{xcolor}
\usepackage{tikz}
% ulem, not soul. `soul`'s \st and \ul are built on a character-by-character
% scan that XeLaTeX's font machinery breaks: the document fails with an
% undefined control sequence rather than degrading. `normalem` stops it
% hijacking \emph, which the transcriptions use for underlining of Jo Hopper's
% own.
\usepackage[normalem]{ulem}
% ulem sets each underlined word in a box TeX can neither hyphenate nor break
% after, so a paragraph dense in \uncertain{} readings can reach a state where
% no legal breakpoint exists. \emergencystretch lets TeX loosen it on a third
% pass instead of overrunning: an overfull line in a transcription hides
% content.
\setlength{\emergencystretch}{3em}
\usepackage{hyperref}
\hypersetup{colorlinks=true,linkcolor=black,urlcolor=black,pdfborder={0 0 0}}

% --- Batch metadata ---------------------------------------------------------
% Taken as they stand by the rendering script, which shows them at the head of
% the reading view. They fix what the file claims to cover: without them a
% transcription floats unmoored in an archive of five hundred sheets.

\newcommand{\ledger}[1]{\gdef\hopLedger{#1}}
\newcommand{\ledgertitle}[1]{\gdef\hopLedgerTitle{#1}}
% A notebook of Josephine Hopper's, in place of a ledger: one file to a
% notebook, filed under transcripts/notebooks/, the argument the site's slug
% (« battle-of-wash-sq »). Sets the ledger to « notebooks » for everything
% that reads it, so the title block and the watermark behave as they do for a
% batch; the rendering checks the sheets against the notebook's own listing.
\newcommand{\notebook}[1]{\gdef\hopLedger{notebooks}\gdef\hopNotebook{#1}}
\gdef\hopNotebook{}
% The Whitney's accession number for the physical volume — 96.208 … 96.213.
% This is what a citation of the object cites.
\newcommand{\objectnumber}[1]{\gdef\hopObject{#1}}
% The batch this file covers. Leaving it out drops the « (batch N) » from the
% title block rather than printing a number that would be false.
\newcommand{\batch}[1]{\gdef\hopBatch{#1}}
\newcommand{\sheets}[2]{\gdef\hopFirst{#1}\gdef\hopLast{#2}}

% The Whitney's date range for the volume, copied verbatim from its resource
% record — never inferred here, never narrowed to the years this batch
% happens to cover.
\newcommand{\dating}[1]{\gdef\hopDating{#1}}

% The legal notice, printed in the title block and repeated as a diagonal
% watermark on every page of the PDF.
%
% This is not boilerplate and it is not optional. The ledgers are © Heirs of
% Josephine N. Hopper, licensed by the Artists Rights Society; the Whitney
% records the object rights as transferred to the Museum. A transcription
% reproduces Josephine Hopper's text. Every file must therefore say what it is
% on its own face, not only in a README that does not travel with the PDF.
% A file without this line gets no watermark, so omitting it is visible in the
% output rather than silent.
\newcommand{\watermark}[1]{\gdef\hopWatermark{#1}}

\gdef\hopLedger{?}\gdef\hopLedgerTitle{}\gdef\hopObject{}\gdef\hopBatch{}%
\gdef\hopFirst{?}\gdef\hopLast{?}\gdef\hopDating{}\gdef\hopWatermark{}

% --- Keywords ---------------------------------------------------------------
%
% Closing a transcription: the vocabulary under which to search for what these
% leaves record — dealers, exhibiting societies, media, places.
% `npm run manifest` extracts them to tag the whole ledger. This line is the
% single source of the tags; there is no tags file, so no tag can describe
% leaves nobody has read.
%
% It is the one thing in the file that is not on the paper, and it earns its
% place: a reader looking for Keppel needs some way in, and a tag written by
% whoever read the sheets is the only kind that cannot describe unread ones.
\newcommand{\keywords}[1]{\par\smallskip\noindent
  {\footnotesize\textsc{Keywords} --- \textit{#1}}\par}

% --- The anchor -------------------------------------------------------------

% The sheet begins here. Two arguments, and both are needed.
%
% #1 is the Whitney's **resource ref** — 16853 — which is the sheet's whole
% address: it is what the image URL is built from and what turns the facsimile
% as the transcript is scrolled. #2 is the number **written on the paper**, or
% empty where nobody wrote one, which is the case for every cover, flyleaf,
% index and loose insertion — about one sheet in nine.
%
% Keeping both is the whole of the numbering problem in this archive. Three
% sequences overlap: the Hoppers' own leaf numbers, the Whitney's order of
% photographs, and the resource refs, which are in no order at all. Only the
% first is citable and only the third is addressable, so a transcription that
% recorded one of them would be either uncitable or unopenable.
%
% `scripts/render.mjs` checks #2 against what `scripts/catalogue.mjs` parsed
% out of the Whitney's descriptor and fails on a disagreement, because a
% mistyped leaf number silently moves an entry to another year.
\newcommand{\sheet}[2]{%
  \marginnote{\footnotesize\color{gray!70}\textsf{\ifx\relax#2\relax---\else#2\fi}}%
  \label{sheet:#1}\ignorespaces}

% --- The critical apparatus -------------------------------------------------

% Illegible: never guessed. On a ledger leaf this is most often a figure, and
% a guessed figure is a fabricated transaction.
\newcommand{\ill}{\textcolor{red!60!black}{[\,\ldots\,]}}

% A doubtful reading: the word or figure is given, and flagged as uncertain.
%
% Underlining belongs to this macro and to nothing else. Jo Hopper underlines
% her own headings --- « Etchings », « Etching Plates », « Early oil » --- and
% those are set with \emph{} instead, because a reader must be able to tell a
% doubtful reading from her emphasis at a glance and the two cannot both be
% underlines. Which of them keeps the underline is decided by consequence: an
% emphasis shown as italic loses nothing, and a doubtful figure shown as a
% certain one loses everything.
\newcommand{\uncertain}[1]{\uline{#1}}

% An editorial addition — an expanded abbreviation, an implied dollar sign.
\newcommand{\add}[1]{\textcolor{blue!50!black}{[#1]}}

% Struck out in the book. Jo Hopper struck a great deal: a price revised, a
% buyer who withdrew, an exhibition that fell through. What she crossed out is
% part of the record and often the more interesting half of it.
\newcommand{\struck}[1]{\sout{#1}}

% The ink it is written in, where the ink says something.
%
% Book IV is a state machine and the colour is its status field: charges in
% black, receipts in red, the sums in pencil, a rule that holds from leaf 7 to
% the end of the volume. A transcription that records the words and the
% figures and nothing about their colour throws that field away, and the
% accounts then have to guess a receipt from the word « rec'd » --- which the
% volume stops writing on leaf 157. So the ink is recorded, per cell, as what
% is on the paper and not as what it means: \ink{red}{rec'd by check} says
% the words are red, and it is `scripts/accounts.mjs` that says red is money
% received. Black is the default and is not marked. The inks are the three
% the volume uses besides black; a fourth would be added here and in
% `scripts/render.mjs` and `scripts/tei.mjs` in the same commit.
% The file is \input, not \usepackage'd, so @ is not a letter here without saying so.
\makeatletter
\@namedef{hop@ink@red}{red!70!black}
\@namedef{hop@ink@pencil}{gray!55!black}
\@namedef{hop@ink@blue}{blue!55!black}
\newcommand{\ink}[2]{%
  \ifcsname hop@ink@#1\endcsname
    \textcolor{\csname hop@ink@#1\endcsname}{#2}%
  \else
    \PackageError{hopper}{\string\ink{#1}: unknown ink}{The inks are red, pencil and blue.}%
  \fi}
\makeatother

% A rule drawn across the money column above this row's figure.
%
% The writer rules off a run of items and writes their sum under the rule,
% and the rule is the whole of what says « this figure is a sum ». It stands
% at the head of the row the sum is on: \ruledoff & & 4355 & 00. On paper it
% is a line under the last two columns of the four a Book IV leaf has; if
% another volume ever rules a different column, the macro grows an argument.
\newcommand{\ruledoff}{\cline{3-4}}

% A diary entry begins. #1 is the date as she wrote it — « Mar. 12 »,
% « Sunday » — and #2 the date the transcriber assigns it, as a year, a month
% and a day (1947-03-12), or as much of one as the leaf allows. The first is
% what is on the paper; the second is the one editorial claim a diary needs
% that a ledger does not, because a ledger's leaf is its own address and an
% entry's date is its subject. Both are kept, apart, so the claim can be
% argued with. Used only in the notebooks; a ledger has no entries.
\newcommand{\entry}[2]{\par\medskip\noindent
  {\bfseries #1}\ifx\relax#2\relax\else\hspace{0.6em}{\footnotesize\color{gray!60!black}\textsf{#2}}\fi
  \par\nopagebreak\smallskip}

% The transcriber's note. Ours, not theirs.
\newcommand{\note}[1]{\par\smallskip{\small\color{gray!60!black}\textit{[#1]}}\par\smallskip}

% A marginal note **of theirs**, distinct from the body — and distinct from
% \note{} for exactly the reason \note{} exists.
\newcommand{\marginal}[1]{\par\smallskip\noindent\hspace{1em}%
  {\small\color{gray!60!black}#1}\par\smallskip}

% --- What is particular to this archive -------------------------------------

% Whose hand wrote it.
%
% This macro has no counterpart in the Grothendieck preamble, and it is the
% single most important thing in this one. These books were written by two
% people and annotated by more. Edward Hopper drew the work and wrote its
% title and size; Josephine Nivison Hopper wrote everything else — the
% descriptions, the anecdotes about people he refused to discuss, the prices,
% the buyers, the columns; and later hands, curatorial and unidentified, added
% to the leaves after 1967.
%
% A transcription that flattens the three has destroyed the archive's chief
% documentary value, which is that it records *two* working lives and the
% division of labour between them. #1 is `edward`, `jo`, `later` or
% `unidentified`; `unidentified` is the right answer far more often than a
% transcriber's confidence suggests, and it is not a failure to use it.
\newcommand{\hand}[2]{%
  \ifx\relax#1\relax#2\else
    \textcolor{gray!55!black}{\footnotesize\textsf{[#1]}}\,#2%
  \fi}

% Edward Hopper's ink record sketch stands here.
%
% The argument is **what is inscribed on or beside the sketch** — a title, a
% dimension, a note — and never a description of the image. Describing the
% drawing would be writing a new document: the sketch is on the site, one pane
% away, at the resolution the Whitney publishes, and prose about it competes
% with looking at it. Leave the argument empty where nothing is written.
% The mark is drawn rather than set from a symbol font: there is no
% mathematics anywhere in this archive, so no maths package is loaded, and
% $\square$ would be an undefined control sequence — which is exactly how it
% failed the first time the PDFs were compiled.
\newcommand{\sketchmark}{\raisebox{0.15ex}{\color{gray!50!black}%
  \tikz[baseline]{\draw[line width=0.4pt] (0,0) rectangle (1.1ex,1.1ex);}}}
\newcommand{\sketch}[1]{\par\smallskip\noindent
  {\small\color{gray!50!black}\sketchmark~\textsf{ink sketch}%
   \ifx\relax#1\relax\else\ --- #1\fi}\par\smallskip}

% Something printed and pasted to the leaf — a reproduction cut from a
% catalogue, a review, a dealer's label. The argument transcribes its printed
% text. These are the only places in the archive where the words are nobody's
% handwriting, and they date the leaf, which handwriting does not.
\newcommand{\clipping}[1]{\par\smallskip\noindent
  {\small\color{gray!50!black}\textsf{[clipping]}\ #1}\par\smallskip}

% A work's block opens here, under the title **exactly as the leaf gives it** —
% « Les Deux Pigeon » and « Les Poillus » stay misspelled, because the
% misspelling is Hopper's and is how the entry is found.
\newcommand{\work}[1]{\par\medskip\noindent{\bfseries #1}\par\nopagebreak\smallskip}

% --- The ruled table --------------------------------------------------------
%
% Jo Hopper ruled these columns herself and kept them for forty years; the
% table is not a presentation of the content, it *is* the content. So it is
% transcribed as a table and rendered as one, on screen and on paper, from one
% source.
%
% #1 is the column specification; #2 is the header row, `&`-separated like any
% other. A cell she left blank is left blank; a cell that cannot be read takes
% \ill{}.
%
% **Use `Y{...}` for any column that can hold a sentence** — an exhibition, a
% dealer's terms, a note — and `l`, `r` or `c` only for dates, figures and
% short names. See the note on \newcolumntype{Y} above for why that is not a
% matter of taste.
%
% The header row is emitted **bare**, between rules, and is deliberately not
% wrapped in \textbf{}. It cannot be: an `&` inside a braced group is not an
% alignment tab, so `\textbf{Date & Exhibitions}` fails at the first header
% with « Forbidden control sequence found while scanning use of
% \check@nocorr@ » — which is TeX saying, at some distance, that the tabs are
% in the wrong place. Rules carry the header instead, which is how a scholarly
% table sets one anyway. (The reading view has no such constraint and renders
% the same row as <th>.)
\newenvironment{ledgertable}[2]
  {\par\smallskip\small
   \begin{longtable}{#1}
   \hline
   #2 \\
   \hline
   \endhead}
  {\end{longtable}\par\smallskip}

% --- The appendix of notes --------------------------------------------------
%
% Everything above the appendix is on the paper. Nothing in it is.
%
% A note says what a named source says about a work these leaves record --- what
% the holding museum measures the canvas at, which other leaf carries the other
% half of a transaction --- and it is the only prose in the file that a reader
% cannot check against the photograph. On the site that is handled by putting
% the note beside a link to its source; on paper a link is not enough, because
% a PDF outlives the site it was downloaded from and travels away from it. So
% every claim here prints its source in full, name and address both, and every
% note prints the day it was written and what wrote it.
%
% The appendix is assembled by `scripts/pdf.mjs` from
% `src/content/work-notes.json` at compile time and is in no `.tex` file. That
% is the point of it. The transcription is the edition; a note is a later and
% weaker kind of writing *about* the edition, and keeping the two in one source
% file would eventually put a sentence somebody inferred inside the document
% whose whole claim on a reader is that nothing in it was inferred.
%
% Only the notes on works this batch's leaves name are printed, so the appendix
% follows the batch rather than repeating the archive, and a batch that has
% earned no notes yet gets no appendix at all --- which is a true statement
% about how much has been checked, and the same statement `work-notes.json`
% makes by being nearly empty.

% A URL is the one string in this archive that can run a hundred characters
% without a space in it, and an overfull box fails the build here because in a
% transcription an overfull box hides content. So the appendix tells TeX where
% a URL may break and sets its apparatus ragged right, rather than leaving the
% line to overrun in a file nobody recompiles.
\def\UrlBreaks{\do\/\do\-\do\.\do\_\do\?\do\&\do\=\do\+}
\Urlmuskip=0mu plus 1mu

\newcommand{\notesappendix}{\clearpage
  \par\noindent{\large\bfseries Notes on the works}\par\smallskip
  \noindent{\footnotesize\color{gray!60!black}\raggedright
    Not on the paper, and not part of the transcription. Each note below
    concerns a work named on a leaf of this batch and says what a named source
    says about it --- a museum's own catalogue record, or another leaf of these
    ledgers. Every sentence carries the source that states it, printed with its
    address so that it can be followed from a copy of this file alone. Where a
    ledger and a museum disagree the disagreement is printed rather than
    resolved.\par}\par\medskip}

% One work: the title as the leaves give it, then the note.
\newcommand{\worknote}[2]{\par\medskip\noindent{\bfseries #1}\par\nopagebreak
  \smallskip\noindent #2\par}

% One claim under it: what is asserted, and who asserts it.
\newcommand{\workclaim}[2]{\par\smallskip\noindent
  {\footnotesize\color{gray!60!black}\raggedright\hangindent=1.6em\hangafter=0
   \hspace{1em}#1\ --- #2\par}}

% Who wrote the note and when, closing the block. A note is a model's reading
% of retrieved records and is weaker evidence than anything above it; the file
% records that rather than letting the appendix pass for the edition.
\newcommand{\worknoteby}[2]{\par\smallskip\noindent
  {\scriptsize\color{gray!55!black}\hspace{1em}Written #1 by #2.\par}\par\smallskip}

% --- Title ------------------------------------------------------------------

% The watermark is drawn on the shipout background so it sits under the text of
% every page, diagonal and light enough to leave the record readable.
\AddToHook{shipout/background}{%
  \ifx\hopWatermark\empty\else
    \begin{tikzpicture}[remember picture, overlay]
      \node[rotate=52, scale=3.4, align=center, text=gray!18,
            font=\sffamily\bfseries]
        at (current page.center) {\hopWatermark};
    \end{tikzpicture}%
  \fi}

\AtBeginDocument{%
  \begin{center}
    {\large\bfseries Edward and Josephine Hopper --- \hopLedgerTitle}\\[2pt]
    \ifx\hopObject\empty\else
      {\small Whitney Museum of American Art, \hopObject}\\[4pt]
    \fi
    {\small sheets \hopFirst--\hopLast
      \ifx\hopBatch\empty\else\ (batch \hopBatch)\fi}\\[2pt]
    \ifx\hopDating\empty\else
      {\small\color{gray!60!black}The Whitney's dating: \hopDating}\\[2pt]
    \fi
    \ifx\hopWatermark\empty\else
      {\small\bfseries\color{red!45!gray}\hopWatermark}\\[2pt]
    \fi
    {\footnotesize\color{gray!60!black}Transcribed from the digitised sheets published by the
     Whitney Museum of American Art.\\
     \textcopyright{} Heirs of Josephine N. Hopper, licensed by Artists Rights Society (ARS),
     New York.\\
     Uncertain readings are underlined, illegible passages marked [\,\ldots\,],
     editorial additions bracketed.}
  \end{center}
  \vspace{1em}}

\endinput


\notebook{three-wash-sq}
\ledgertitle{Josephine Hopper --- Re: 3 Wash Sq}
\sheets{1}{51}
\dating{1947}
\watermark{Demonstration edition}

\begin{document}

\section{Front cover}

\sheet{96069}{}
\hand{jo}{Re. 3 Wash. Sq}

\note{in pencil, written across the stationer's printed cover, over the panel
above the « Gripper.edge » device, with a rough line drawn round it and a
flourish under it; the mark after « Wash » is a small open hook that may be a
point or a dash}

\section{Page 1}

\sheet{97207}{1}
\hand{jo}{Interested in Vil. development}

\note{stands above the first ruled line, at the head of the page; the whole of
this page is in pencil}

\hand{jo}{Wash Sq. As. \\
Gr. Vil' As. \\
5 Ave. " \\
Wash. Sq. Neigh' \\
Village Little Gardens \\
Village Art Center \\
\uncertain{Yoan} Village - \uncertain{past.}}

\hand{jo}{Are people more \uncertain{sympathetic} \\
than property - protection \\
1" step zoning - \\
uniform restrictions}

\note{the seven names are gathered by a brace, and the four lines beginning
« Are people » stand to the right of it}

\hand{jo}{S.S. \uncertain{H.} Trustees \\
Pres - N.Y. State Chamber Commerce \\
\uncertain{57} Liberty \qquad Chairman of Bd - \\
Peter Grimm \\
Mayor of N.Y. ex officio \\
Trinity - Episcopal Ch. \\
12" St + 5" \quad Presbyterian Ch. \\
2 Seamans assoc.}

\hand{jo}{officers \\
\uncertain{Sailors}}

\note{« S.S. H. Trustees » is underlined with a stroke that sweeps up to the
right; « Peter Grimm » is in a lighter pencil than the lines around it; the
last two lines are gathered by a brace standing beside « 2 Seamans assoc. »}

\hand{jo}{Supposed to be a \\
non profit organization}

\note{stands to the right of « Mayor of N.Y. ex officio » and the two lines
under it; « Supposed to be a » is underlined and « profit » is written over}

\hand{jo}{Comptroller \ill{} \qquad Impersonal board.}

\note{the word after « Comptroller » is written twice over itself; the first
letter is an F or an h and the ending will not be read}

\hand{jo}{La Guardia - WJZ - noon Sun. for Kraft Cheese - Steinberg \\
30 Rockefeller Plaza - 50" floor. Miss Steinberg. \\
Residence: Bronxville, N.Y. \qquad Hon. F.H. LaGuardia ? \quad Secretary}

\note{« Steinberg » stands at the fore-edge of the first line and again on the
second, where it is broken across two lines; a query mark of hers stands below
the last line between the name and « Secretary », and a stroke runs over
« Secretary »}

\hand{jo}{Wash Sq. As.- Pres. King \uncertain{Woodredge} \\
Vice " - \uncertain{J.} Stanley Hillyer - Chairman of Zoning Committee \\
Sec. - Ed. Steinhart.}

\hand{jo}{Gilbert Hodges - Ex Gen. manager of Sun}

\note{stands to the right of the three officers; « Committee » is written down
the fore-edge, the pad turned. Battle of Wash Sq. gives the same man as
« H. Stanley Hillyer »; the initial here is a J or an I and the stroke will not
settle it}

\section{Page 2}

\sheet{96239}{2}
\note{this photograph, and every one that follows in this sitting, shows two
written sides of the pad: the leaf turned back above the spiral, and the leaf
below it. They run on continuously and are transcribed in that order}

\hand{jo}{If Mr. \uncertain{Cortisoz} were still on deck, he would \\
be really concerned to learn what \\
is happening in the \struck{country} the art \\
world of the Nation. Surely Wash. Sq. \\
\struck{has} is entitled to that \struck{distinction} as it is as \\
\struck{on} the Latin Quarter is to France. \\
Wash. Sq. \struck{\ill{}} surrounds a famous \\
park, of which the N.Y.U. has acquired all of the E. side \\
+ various other holdings on the S. side. The truth \\
now emerging that the N.Y.U. proposes to \\
obtain possession of all 4 sides, \struck{\ill{}} watching \\
\struck{property} \ill{} swallowed up so as to have the Park as its Campus. \\
It managed to gobble up the old Judson + several studio \\
buildings to the West, not hesitating to \\
oust John Sloan from \struck{the} his studio he \\
had occupied \struck{for} many years * . This deed \\
created small outcry. But a local clamor \\
\ill{} is growing \struck{when + the} from the event of N.Y.U. \struck{has} securing a \\
21 year lease from the S.S.H. \struck{\ill{}} \\
the famous old \# 3, for great many yrs. \struck{generations} the \\
stronghold of generations of artists. \struck{Its} \\
record of \# 3 is historic - + includes such names among \\
former \struck{tenants} tenants as - Wm. Glackens + his wife Edith \\
Dimock}

\marginal{still \quad --- above the line, with a caret after « were »}

\marginal{Even \quad --- before « as it is as », over the struck
« distinction »}

\marginal{that \ill{} of green + trees \quad --- above the struck
« property »}

\marginal{his \quad --- written over the struck « the »}

\marginal{J.S. had \struck{\ill{}} immortalized this region of his residence,
the \struck{\ill{}} Square in its various aspects + activities, street scenes,
also the \struck{\ill{}} elevated, Jefferson Market, on west to the ferry boats
plowing thru heavy atmosphere, \struck{changing} \ill{} a sojourn in water ways.
This studio was fertile soil for Amer. Art history. \ill{} with proceeds to
\uncertain{evict} its tenants, \ill{} more ready to resist ousting than John
Sloan, grown old in waging battles in behalf of fellow artists. He
\struck{was} one of the founders of the Independent in that far back day. He
taught a generation at the Art Students League + fought mightily in any
\struck{cause} cause affecting them. \quad --- the note the asterisk after
« many years » points at, written in six lines up the fore-edge with the pad
turned}

\note{the first draft, written fast and corrected over itself in the same
pencil; the interlineations are given where they can be read and the passage is
not everywhere recoverable. Page 4 writes the critic's name plainly as
« Royal Cortissoz », which settles who is meant without settling this stroke.
The marginal up the fore-edge is crowded and its middle is partly unread}

\hand{jo}{Mrs. Guy Pène du Bois, Ernest Lawson, Rockwell \\
Kent, Alexander Brook, Rico LeBrun, \uncertain{Conceta} \\
\uncertain{Scaravaglioni}, Coulton Waugh, Marjorie \uncertain{Ryerson} \\
many other painters - + among writers, John Dos \\
Passos wrote his Three Soldiers here, Schofield \\
Thayer was a resident while publishing the \\
ultra sophisticated Dial, E.E. Cummings, \\
Edmund Wilson, Eleanor Wylie, Jack Kelley \\
who with S. Thayer + some others sailed a tiny boat \\
across the Atlantic. To go back farther, there \\
were Frank Harris holding his famous \\
Salon, Mary Tillinghast designing her \\
stained glass + entertaining the young \\
Paderewski, who gave an impromptu recital \\
in her studio; Benj. Porter over 35 yrs. ago \\
painting stylish portraits, Grace \uncertain{Randolph} \\
founding the Nat. As. of Women Painters + Sculptors \\
+ Thos. Eakins using the studio of Fred. Stokes \\
to paint one of his now famous portraits \\
Many other painters + others the last enlivened by \\
L. + D. Gish + Jack Kelley who with Schofield Thayer + some other wild spirits managed \\
to cross the Atlantic in some tiny shell of \\
a boat. Lastly in past + present membership one finds 9 known members of the \uncertain{Inst.} Arts \\
+ Letters}

\marginal{Waldo Pierce \quad --- above « Alexander Brook »}

\marginal{on his 1st trip to America \quad --- above « Paderewski », with a
caret}

\marginal{Mary Blair "Mad" \quad --- above « + Jack Kelley »}

\note{« + Letters » is written below the end of the last line rather than after
it; « Schofield » and « membership » each turn the line, the first with her
hyphen; page 10 settles the first pair of initials as Lillian + Dorothy Gish}

\section{Page 3}

\sheet{96255}{3}
\hand{jo}{no small distinction for one single house \\
Among present residents there are \\
Fred W. Stokes who was artist on staff of \\
Adm. Peary on his trips in search of the \\
N. Pole. Mr. Stokes has lived here nearly \\
1/2 a century + is now 88 - \\
E. Hopper has just turned the 1/3 mark of \\
a century, having spent the greater part of \\
his life \struck{in the} here + painted pictures that hang \\
in museums clean across a continent \\
in this old house: Walter Pach + his wife \\
Magda Pach \struck{reside} had hoped to live + die \\
here; there is also Jos. Pollet, \uncertain{Ganvor Bul} \\
\uncertain{Tielman}, \struck{Louise} whose murals hang in \\
the Times section of Rockefeller Centre; \\
\struck{Louise Moore} Jane Gray doing portraits, \\
Louise Moore, Mrs. Van \uncertain{Meter}, preparing \\
\ill{}, engravers + others. \\
All the hard working productive artists \\
are doomed to eviction by the N.Y.U. \\
whose claim to contributing to Amer. \\
culture is challenged by this deed of \\
vandalism, added to which is the eviction \\
of the unprotesting John Sloan from the Judson, \\
\struck{which act} the memory of which act lives on \\
in art history.}

\marginal{on last day of Jan. \quad --- above « has just turned »}

\note{« Tielman » turns the line after « Bul T »; pages 7 and 10 write the same
name again and no two of the three agree. Page 10 gives « Madame Van Meter »
plainly, which is why the reading is offered here}

\hand{jo}{\# 3 does not take this matter philosophically. \\
If they are creators of art, they must also \\
be creators of sentiment. \struck{they protest} \\
They conduct meetings of protest, tell their \\
tale, write letters, make phone calls, \\
try to enlist the concern of associations \\
If these turn a deaf ear, there is a chance \\
that not only 3 will be swallowed up, \\
but all \struck{too soon}. this power of aggression, \\
\struck{unchecked} unarrested, will be in \uncertain{possession} \\
of a City Park. John Doe, + his \\
wife + baby carriage can go elsewhere. \\
If the S.S.H. wishing to escape \struck{taxes} \\
taxes on the land, lets go the cluster of \\
\struck{so extremely beautiful} the houses with those extremely \\
beautiful facades - incl. in \# 7 thru 13. W.S.N. \\
the Univ. will have almost reached from Univ. Pl. to 5 Ave. \\
John \uncertain{Morron} with his 2 stone lions, \\
(by \struck{\ill{}} Kitson 1883) \uncertain{blandly} valiant \\
guard over \# 6 - may their strength \\
avail - to check the onslaught. \\
The wife of one of the painters, herself a former \\
Claims - "It's all just like Hitler, reaching \\
Vienna. No one stopped him either}

\note{« possession » turns the line and its second half is doubtful}

\section{Page 4}

\sheet{97352}{4}
\hand{jo}{If Mr. Royal Cortissoz were still on deck at the Tribune \\
he would be really concerned \struck{by} to learn of \\
what is happening in a small section of \\
this City that is analogous to the Latin Quarter \\
\struck{known} that has long been Washington Square, is now the center \\
of art in America. Such is the unrivaled prestige of W.S. \\
This Square surrounds a famous Park, which is now \\
\struck{endangered by} exposed to the predatory \struck{ambitions} encroachments of \\
the N.Y.U. The N.Y.U. now occupys all of the E. side \\
+ various holding on the S. side. The danger \struck{to} \\
now makes itself felt that the N.Y.U. \struck{will} wishes to \\
\struck{absorb} gain possession of all 4 sides surrounding \\
+ thus acquire the Park itself for a Campus. \\
It has managed to gobble up the fine old Judson \\
+ several studio buildings to th. S. farther West \\
not hesitating to oust no less an important figure \\
than John Sloan from his fine studio adjoining \\
the Judson. In this studio J.S. created many \\
memorable canvases depicting scenes in this \\
neighborhood, the Park in its various aspects. \\
street scenes, the old 6" Ave. Elevated careening \struck{in its} \\
sweep around a dizzy corner, the old Jefferson \\
Market, + farther West, ferry boats toiling thru \\
City water ways in \struck{thick} dirty weather. These canvases are Historic}

\marginal{in the opinion of many \quad --- above « exposed »}

\marginal{one of these \quad --- above « wishes to », over the struck « will »}

\marginal{on the S. \quad --- above « farther West »}

\marginal{careening \quad --- above the struck « in its »}

\note{the second draft, in a lighter pencil and a slower hand than page 2; a
few of the interlineations at the fore-edge are crowded and are not read}

\hand{jo}{record, told with much vigor + abiding beauty. \\
It will be impossible for future historians \\
to ignore the dastardly abuse of this man's \\
comfort + disturbance of his work by \\
thrusting him forth to find suitable quarters elsewhere. None \struck{were} available in \\
his accustomed neighborhood. It so happens \\
To-day no studios \struck{none} are available in any neighborhood - and a clamor is now \struck{spread} forth when \\
the N.Y.U. secured a 21 year lease from the S.S.H. \\
+ began processes to evict the tenants \\
of \# 3 W.S. North. 3 Wash Sq. N. is one of \\
that row of perfect houses on the N. side of the \\
Sq. just E. of 5 Ave. It has been \struck{is} a famous old house \\
for over a century the home of many \struck{people} \\
\ill{} of \struck{high} distinction in the art world. \\
Among former tenants there is included. \\
Wm. Glackens + his wife Edith Dimock, \\
the Guy Pène duBois', the Rockwell Kents \\
Ernest Lawson, Ricco LeBrun, Alex Brook, \\
Waldo Pierce, \uncertain{Concerta Scaravaglioni}, Marjorie \\
Ryerson, Coulton Waugh + many others. Among the \\
writers are John Dos Passos who wrote his 3 Soldiers \\
in this house, Schofield Thayer publishing his \\
most sophistocated Dial, Edmund Wilson, \\
E.E. Cummings, the lovely poetess Eleanor Wylie}

\marginal{He was practically exiled \quad --- above « It so happens »}

\note{« quarters » and « neighborhood » each turn the line, both with her hyphen}

\section{Page 5}

\sheet{97295}{5}
\hand{jo}{\struck{When there was} \\
\struck{Grace Rando} Ever farther back is Grace \uncertain{Randol}, \\
founding the Nat. As. of Women Painters + \\
Sculptors, Mary Tillinghast designing \\
stained glass + entertaining the young Paderouski on his first trip to America. It is told \\
that he gave an impromptu performance \\
in her studio. Also it delights present \\
tenants to hear that Thos. Eakins painted a \\
portrait on the 4" floor back, using the \\
studio of his friend Fred W. Stokes who at \\
89 is still there. The list of celebrities \\
is enlivened by including \struck{the} L. + D. Gish, \\
+ Mary Blair actress + "Mad" Jack Kelly \\
who with Schofield Thayer + other wild \\
Indians \struck{now} fresh out of \struck{college from} Harvard \\
crossed the Atlantic (in some tiny "shell"). Going even farther \\
back there was F. Hopkinson Smith writing his \\
Tales of No 3 + entertaining friends about an \\
enormous fireplace with lion heads on either side. \\
there was \uncertain{Carol Delvaille} with a canvas in \\
the Luxembourg + Frank Harris holding \\
his famous Salon - others too numerous to \\
list here. These latter are the ghosts whom the \\
\struck{present} tenants, \struck{awaiting} combating eviction would fain \\
evoke to haunt the nefarious Univ. that seeks}

\marginal{he \quad --- after « Jack Kelly », at the end of the line}

\marginal{most perilously \quad --- above « crossed the »}

\marginal{these lions still \struck{stand} hold their ground - \quad --- up the
fore-edge beside the fireplace, the pad turned}

\note{« Paderouski » turns the line}

\hand{jo}{to send them forth \struck{with their experience} with not a studio available. \\
The present tenants, nearly all painters \\
incl: their dean, Fred W. Stokes, who has lived \\
in this house nearly half a century. * E.H. has \\
painted here for \struck{33} exactly 33 1/2 yrs. which \\
makes 1/3 century, his wife a \struck{also a} painter, \struck{has} \\
Rockwell Kent's old studio. Walter Pach + his \\
painter wife Magda Pach, Jos. Pollet with a fine \\
landscape at the recent Whitney show, \uncertain{Ganvor Bul} \\
\uncertain{Tielman} with murals at Rockefeller Center. \\
* who was artist of staff with Adm. \\
Peary on his trips in search of the North Pole. \\
+ who went on a later Expedition with Amundsen. He is now 88 yrs. old + climbs his \\
74 steps \struck{up} like a mountain goat. He has \\
done this for nearly half a century.) \\
Also there is Jane Gray painting portraits, Louise \\
Moore designing beautiful jewelry, Madame \\
Van \uncertain{Motes} \struck{preparing} under her high ceilings \\
preparing operatic song birds. And others, \\
including 2 Colonels - one ex Ziegfeld Folly girl now \\
taking up art, living with the lions on Hopkinson \\
Smith's huge fireplace - + her mother entertaining Rallies for Non Eviction. They all believe \\
it ill befits a university, professing \struck{to}}

\note{the asterisked passage is enclosed by a drawn line and stands in the run
of the prose, not in the margin; « Amundsen » and « entertaining » each turn the
line. The first side
gives Stokes's age as 89 and the second as 88; both stand as written}

\section{Page 6}

\sheet{96070}{6}
\hand{jo}{at least to endorse, Amer. Culture, to evict for \\
their own selfish + ill defined uses 14 \\
\struck{producing} creative artists, all in the \\
act of producing work of recognized value as contributions \\
\struck{to national cultural output} \\
accomplishments. * \struck{paper}}

\entry{Tues. Feb. 19"}{1947-02-19}
\hand{jo}{\struck{they} appear for a \\
hearing before the O.P.A. to state their \\
reasons for Non-Eviction. If their case \\
is lost that much more territory is gained \\
by the N.Y.U. in their Course of Conquest. \\
\struck{So} "Recalling only too well the \struck{march} recent descent of \\
Hitler on Vienna" claims one of the \\
outraged lady artists. "Ta \uncertain{bouche}, bébé - \\
counsels her calmer husband. "Wait + See." \\
* "In an earlier age, these people would \\
have routed even the Transcendentalists! \\
(to their ever lasting shame) \\
"Well, while we are about it - why not \\
page Euripides!" \\
"Scribes + Pharasees (they are) hypocrites!"}

\marginal{painter people \quad --- above the struck « they »}

\marginal{around the Square \quad --- after « is gained »}

\hand{jo}{To their everlasting shame, remembering what they did to John Sloan. In an \\
* \struck{Sure, to their} sure these people in an earlier age would have \\
evicted even the Transcendentalists!"}

\note{written small at the foot of the side, below the ruling that carries the
rest}

\entry{Tues. Feb. 19"}{1947-02-19}
\hand{jo}{\struck{they} appear for a hearing before the \\
O.P.A. to state their reasons for non-eviction. \\
If their case is lost, that much territory is \\
gained by the N.Y.U. in their Course of Conquest \\
around Washington Square. \\
"Recalling only too well the recent descent \\
of Hitler on Vienna" claims one of the \\
outraged lady artists. \\
"Ta boucha, bébé" counsels her calmer \\
husband. "Wait and See." \\
"To their everlasting shame, I remember \\
what they did to John Sloan. In an earlier \\
age, these people would have evicted even \\
the Transcendentalists!" \\
"Well, while we're about it, why not page \\
Euripides?" \\
"Scribes + Pharasees - hypocrites!" \\
\struck{+ may their ears burn to cinders,} \\
May they toss in \struck{in} their "unquiet graves" \\
"Or find no grave vacant to do their tossing in \\
\struck{The eviction of aborigines painters from \# 3 will clear that much more} \\
\struck{park territory for N.Y.U. aggression.} \\
and at that Wash. Sq. was once a graveyard. Now children play there, + the trees are green in Summer. \\
in the heart of a city. It is a delightful neighborhood. Why permit that it should become a College Campus.}

\marginal{A.M \quad --- above the date}

\marginal{these painter folks \quad --- above the struck « they »}

\marginal{in Concord \quad --- above « age these people »}

\marginal{of a city \quad --- above « in the heart »}

\note{the passage is written out twice on this one photograph, first on the
upper side and then on the lower; the last three lines of the lower side are
written smaller and crowded into the foot of the page. « bouche » on the upper
side and « boucha » on the lower are as they stand; page 11 writes « bouche »
in a careful, almost printed hand}

\section{Page 7}

\sheet{97317}{7}
\hand{jo}{\struck{Strictly Confidential} - this preword for explanation \\
Because of the biographical sketch in last week's \\
New Yorker, I'm venturing to try to enlist your \\
interest in the plight of an old studio house \\
\struck{much} a terrible old house with no heat, \struck{no} \\
high ceilings, long flights to climb, bath rooms - \\
remote but \struck{beautiful} even the bare walls beautiful \\
A house frequently cussed but forever adored. \\
Not only the house - but the people who do cherish it \\
+ are being heaved out - by \struck{with} what considerable \\
exaggeration of merit one calls a 3" class University. \\
The chancellor of that institution occupys \\
\# 5 Wash Sq N. - We live in \# 3 - one house \\
separating him from our spleen. His house, \\
like ours is 1/2 a city block in depth - \\
+ he, his wife + son occupy the whole of \struck{\# 5}. But \\
he would throw 14 creative artists + several others \\
out in the snow so as to gain possession of still \\
more park \struck{territory}, during college vacation. \\
(Long years ago, I worked on the old Tribune \\
\struck{for a short time} .) \\
We are told our chances legally are of the slightest \\
but public opinion might deter an institution \\
\struck{with} already given 8 000 000 but making a drive \\
for more, from doing \struck{the} worst - to us. --- at an OPA \\
hearing Tues. A.M. Feb. 19" - or thereafter.}

\marginal{(furnished) \quad --- above « with no heat »}

\marginal{\ill{} \quad --- above « son », two words not read}

\marginal{too \quad --- above « Long years ago »}

\entry{Feb. 10, 47.}{1947-02-10}
\hand{jo}{Mr. Whitelaw Reid - Herald Tribune}

\hand{jo}{Everybody in \# 3 is now scurrying about to \\
stir up anything it can, public opinion \\
in \struck{their} its own behalf. I submit this highly prejudiced \\
version of the news. We are so close to the fact, a \\
sense of fact is acutely realized - yet what I \\
write is no clairvoyance - but substantiated \\
by counting real estate accretions + their direction. \\
Please forgive all this long hand. I shall try to \\
find some one in the neighborhood to type my \\
main story in the morning. \\
If you have struggled thru this, I thank you \\
heartily. Sincerely}

\marginal{notable \quad --- above « their direction »}

\hand{jo}{This part separate = preword \\
Strictly Confidential this preword - please.}

\note{under a rule, below the signature; « at top circled in red » is
interlined in the first of the two lines, and nothing on this side is in red}

\hand{jo}{"May they toss in their unquiet graves" \\
"Or find no graves vacant to do their tossing in." \\
And at that Wash. Sq. was once a grave yard! Now \\
children play + in the spring trees are green, \struck{the place} an oasis \\
a delight to the neighborhood. \struck{Why} In the heart of a densely \\
crowded city why permit \struck{that} ever that it \struck{should} become \\
a college campus. Why permit a gigantic university, \\
\struck{tax free}, to occupy large spaces, \struck{tax} free, in the midst \\
of a great city. Set them do their expanding elsewhere. Who wants \\
them anyhow. They evict tax paying artists aborigines, who trust \\
that now an alert has gone forth. Let all brave men + true}

\marginal{do their expanding \quad --- above « Set them »}

\marginal{paying \quad --- above « tax artists »}

\section{Page 8}

\sheet{96270}{8}
\hand{jo}{\struck{about} from which so much of the \\
what is happening to Wash. Sq. - a small section \\
of the City \struck{that has} about which is gathered so much of \\
the art interests of the nation. \\
Join the embattled artists to save the park + save \\
historic \# 3 for the artists as their birth right. \\
all in the act of contributing to that great fund \\
while defending the traditions of the fine old \\
house + defending it against agression}

\note{the whole of this side is cancelled with a large cross drawn over it;
« as their birth right » is interlined after « artists », and the rest of the
side is blank}

\entry{Feb. 10, 47.}{1947-02-10}
\hand{jo}{Mr. Whitelaw Reid, Herald Tribune, \\
\struck{Hope to get} (this typed) - \\
If Mr. Royal Cortisoz were still on deck at the \\
Herald Tribune he would be genuinely concerned \\
to learn what is happening in Washington Square, \\
a small section of the City about which is gathered \\
so much of the art interests of the nation. \\
Washington Square surrounds a famous park, \\
all of the East side of which the New York University \\
is in possession. In the south side it has also \\
large holdings. Recently it has acquired from the \\
Sailors' Snug Harbor, which is not permitted to \\
sell any of its territory, 21 year leases of the \\
first three houses at the North corner of the Square, \\
consisting of 3 of the original houses of that \\
famous row on Washington Square North to \\
the east of 5" Ave. The danger now makes \\
itself felt that the N.Y.U., unchecked, \\
will gain possession of all 4 surrounding sides \\
\struck{the} and make of the park itself the College Campus. \\
This fear is not substantially unfounded. \\
The N.Y.U. has managed to gobble up the fine \\
old Judson and several studio buildings \\
farther west of the Judson on the S. side of the Sq. \\
over.}

\marginal{east \quad --- above « North corner »}

\note{the fair copy begins here; she numbers its sides 2 to 7 in the head
margin from page 9 onward, and « over. » stands alone under the last line.
« important » and « historians » each turn the line}

\section{Page 9}

\sheet{96242}{9}
\hand{jo}{2}

\hand{jo}{not hesitating to oust no less an important figure than John Sloan from his fine \\
studio adjoining the Judson. In this studio, \\
with the largest skylight + north window on the Sq. \\
John Sloan created many memorable canvases \\
depicting scenes in this neighborhood: \\
the Park in its various aspects, street scenes, \\
the old 6" Ave. elevated careening around \\
dizzy corners, the old Jefferson Market building \\
and farther west, ferry boats toiling thru City \\
water ways in dirty weather. These canvases \\
are historic records painted with much \\
vigor + abiding beauty. John Sloan was \\
not a young man when this eviction took \\
place. It will be impossible for future historians to ignore the dastardly disregard for \\
the comfort of this ageing man and \\
the interruption of his important, valuable \\
work, thrusting him forth to find appropriate \\
quarters elsewhere with naturally none \\
available in his neighborhood. Good studios \\
are rare indeed \struck{at} any time. Real Estate \\
people are concerned with artists only}

\marginal{lively \quad --- above « street scenes »}

\hand{jo}{3}

\hand{jo}{in that with the spread of reputations of the artists \\
\struck{the} real estate shares, tourists arrive, \\
rents go up, the old story, then the artists \\
must go; make room to accomodate \struck{expensive} \\
renovations - or other interests arrive. John Sloan was \struck{thus} \\
practically exiled to find refuge in Chelsea Village, \\
from which appear no more scenes of Greenwich \\
Village and that University continues to grab. \\
To-day no studios are available in any \\
neighborhood. A clamor now goes forth \\
when the N.Y.U. starts procedures to evict \\
the tenants of 3 Washington Square North. \\
This house is one of that perfect row and a \\
very famous old studio building, for over \\
a century the home of many people distinguished in the art world. Former tenants \\
\struck{there is} include William Glackens + his wife \\
Edith Dimmock, the Guy Pène du Bois', \\
Ernest Lawson, Rockwell Kent, Ricco LeBrun, \\
Alexander Brook, Waldo Pierce, Coulton Waugh, \\
Concerta Scaravaglioni, Marjorie Ryerson, \\
+ many others. Among the writers are \\
John Dos Passos who wrote his Three Soldiers \\
here, Schofield Thayer publishing his most \\
sophisticated Dial, Edmund Wilson,}

\marginal{with \quad --- above « in that »}

\marginal{that also \quad --- above « interests arrive »}

\note{« 2 » and « 3 » stand alone in the head margin of the two sides: her own
pagination of the fair copy, not the pad's; « distinguished » turns the line}

\section{Page 10}

\sheet{97561}{10}
\hand{jo}{4}

\hand{jo}{E.E. Cummings; the lovely poetess, Eleanor \\
Wylie. Farther back: Grace Randolph, founding \\
the National Association of Women Painters + \\
Sculptors, Mary Tillinghast designing her \\
famous stained glass + entertaining the young \\
Paderouski on his first trip to America. \\
It is told that he gave an impromptu \\
recital in her studio on the first floor. \\
Also it delights present tenants to learn \\
that Thomas Eakins painted a portrait on \\
the 4" floor rear using the studio of his \\
friend Frederick W. Stokes who at 88 is \\
still there. The list of celebrities is enlivened \\
by including Lillian + Dorothy Gish, + \\
Mary Blair, actress; + "Mad" Jack Kelley \\
who with Schofield Thayer and a gang of wild \\
Indians fresh out of Harvard in some tiny \\
shell of a boat most perillously crossed the \\
Atlantic. Going even farther back, there was \\
F. Hopkinson Smith writing his "Tales of \\
No. 3" and entertaining friends about an \\
enormous fireplace with heads of lions on \\
either side. The lions still hold their ground. \\
There was Carol Delvaille with a canvas in the \\
Luxembourg + Frank Harris holding}

\hand{jo}{5}

\hand{jo}{his famous salon with a handsome big blond \\
wife and \struck{others} \struck{too} numerous to list here now. \\
The present tenants, combatting eviction \\
are nearly all painters; their dean, \\
Frederick W. Stokes who was artist of staff \\
of Admiral Peary on his trips \struck{in search} \\
to the North Pole. He also went on later arctic \\
expeditions with Amundsen. He is now \\
88 years old and until recently, before the \\
shadow of eviction appeared to lay him low, \\
climbed up his 74 steps like a mountain \\
goat. He has done this leg strengthening \\
exercise to his 4" floor at No. 3 for nearly \\
half a century. Other stair climbing painters \\
are Mr. + Mrs. Edward Hopper, he, longer at it \\
than she, on Jan. 31" last, finished his 1/3 \\
century on the top floor of that old house. \\
Other painters: Mr. + Mrs. Walter Pach also \\
climb. Joseph Pollet with a fine painting at \\
a recent Whitney climbs most of all to a really \\
splendid studio. \uncertain{Ganvor Bul Tielman} has \\
murals at Rockefeller Center. Also there is \\
Jane Gray painting portraits, Louisa Moore \\
designing unique and beautiful jewelry, \\
Madame Van Meter under her high ceilings,}

\marginal{people \quad --- above « others too »}

\marginal{on a 5" floor \quad --- above « splendid studio »}

\note{« 4 » and « 5 » stand alone in the head margin; a strip of adhesive tape
crosses the foot of the lower side without covering a word}

\section{Page 11}

\sheet{97079}{11}
\hand{jo}{6}

\hand{jo}{preparing operatic song birds. Others \\
included are two present world war Colonels, \\
\struck{the wife of one of them now painting local scenes +} \\
and living with \struck{her mother} and the lions on \\
F. Hopkinson's old fireplace and entertaining \\
Non Eviction rallies on the second floor rear. \\
They all believe it ill befits a University, \\
professing, at least, to endorse American \\
culture to evict for their own selfish + \\
ill defined uses 14 creative artists, all \\
in the act of contributing to that great fund \\
while maintaining the tradition of the old \\
house so greatly cherished and defending it \\
from agression.}

\note{a word between the eighth and ninth lines is blotted out in heavy black
and cannot be read}

\entry{Mon. A.M. Feb. 18"}{1947-02-18}
\hand{jo}{these painter folks appear for \\
a hearing before the O.P.A. to state their reasons \\
for non-eviction. If their case is lost, (the \\
weight of public opinion can help them) that much \\
territory is gained by the N.Y.U. in their Course \\
of Conquest around Washington Sq. \\
"Recalling only too well the \uncertain{resent} descent \\
of Hitler on Vienna where no one stopped him," \\
claims one of the outraged lady artists. \\
"Ta bouche, bébé," counseled her}

\hand{jo}{7}

\hand{jo}{calmer husband. "Wait and see." \\
"To their everlasting shame I remember \\
what they did to John Sloan. Way back in \\
Concord, that gang would have evicted \\
the Transcendentalists." \\
"Well while we're about it why not page \\
Euripides?" \\
"Scribes + Pharasees --- hypocrites! \\
"May they early toss in their "unquiet graves!" \\
"Or find no grave vacant to do their \\
tossing in!" \\
And as to that, Wash. Sq. was once a graveyard! \\
Now children play there + in the spring the trees \\
are green, the place an oasis, a delight to the \\
neighborhood. In the heart of a densely crowded city \\
why ever permit this become a college campus? \\
Why allow a gigantic university to occupy \\
large spaces tax free in the midst of a great city? \\
Let them do their expanding elsewhere. Who wants \\
them here anyway? They evict taxpaying artist aborigines. \struck{who trust that} Now an alerte has gone forth, \\
let all good men + true join the embattled artists to save \\
the park and also historic No. 3 for the artists as \\
their birth right.}

\marginal{park to \quad --- above « permit this become », with a caret}

\marginal{that \quad --- above « Now an alerte »}

\note{« 6 » and « 7 » stand alone in the head margin; « aborigines » turns the
line; the fair copy ends here, and the sitting with it}

\section{Page 12}

\sheet{97449}{12}
\hand{jo}{It has been felt that the Sq. belongs \struck{with} John Doe + family \\
\struck{belongs to the artists} who \struck{shares it}; + the Sq. to the \\
artists + those concerned with his \struck{\ill{}} output; art + \\
its preservation + wellbeing.}

\marginal{roughly stated \quad --- above « felt that »}

\marginal{Park itself to \quad --- above « the Sq. belongs with »}

\hand{jo}{It has been felt, roughly stated, that the Park \\
belongs to John Doe + family + the Sq. to the artists. \\
The N.Y.U. is \uncertain{poaching} on the artists (+ eventually \\
\#3 protests eviction +).}

\marginal{John Doe also \quad --- at the foot of the passage, to the right}

\note{the first paragraph is the second written out and then corrected over
itself; « poaching » is written over a struck word and neither is certain. The
upper side ends here and the rest of it is blank}

\hand{jo}{Ernest O. Melby - Dean Sch. of Ed. \\
Press Bldg. \struck{\ill{}} Wash. Place. \\
Walter \uncertain{Cook} - lecturer on Hist. of Art. \\
\struck{Rasch} Raasch - Gen. Manager \\
\qquad N.Y.U. \\
Bd. of Directors of N.Y.U.}

\marginal{office' \quad --- above « Melby »}

\marginal{Sp. 7-2000. \quad --- at the fore-edge of the « Melby » line}

\note{« Wash. Place » is underlined; the lower side of the photograph begins
here and is a page of notes, not of draft, in the same pencil}

\hand{jo}{Village \uncertain{Sett} Shop \\
1 Charles St.}

\note{drawn round with a box}

\hand{jo}{Alfred Paul Waldheim \\
partner of Carl Hart.}

\note{stands to the right of the boxed item, with an asterisk above it}

\hand{jo}{O.P.A. \qquad Concerned over residential \\
State Laws \qquad " \qquad business \\
\qquad assessment \\
150 000 tax assessment \#1 \\
60 000 are \quad " \qquad \# 9 no frontage since 1942. \\
\qquad no meter change. \\
80 000 - 22.88 . \\
\qquad + frontage \\
\#4 \quad 70 000 - \\
\#5 - 80 000 - no frontage. -}

\hand{jo}{\uncertain{76/3} \\
Mortage \uncertain{Sifs} Due}

\note{stands to the right of the « 80 000 - 22.88 » line}

\hand{jo}{\#4 \struck{\ill{}} 6 leases . \qquad 1939 starting}

\note{a lone « \# » and a stray looped mark stand between this line and the
one above it; the whole of this side is in pencil}

\section{Page 13}

\sheet{96271}{13}
\hand{jo}{\uncertain{Mrs.} of \struck{\ill{}} support tenants. At next meeting \\
Ctt. will propose movement}

\marginal{receptions of \quad --- above « of \ill{} »}

\note{a long stroke crosses the first line at x-height; « movement » is
underlined}

\hand{jo}{6 1/2 \qquad 300 \\
\qquad \qquad 25 \\
\qquad \qquad 325.}

\hand{jo}{Farrell N.Y. Times \\
assoc. Press relief ?}

\note{stands to the right of the three figures; « Farrell N.Y. Times » is
underlined}

\hand{jo}{Vice Commander, Legion \uncertain{Carett.}}

\marginal{Amer. \quad --- above « Legion »}

\hand{jo}{Newman Cl. - \\
W. M. C. A. \\
(Hebrew organiz.)}

\hand{jo}{on S. of Wash. Sq. \\
in \# 2 Wash. Sq. N.}

\note{a brace gathers the three organisations and the two lines beginning
« on S. » stand to the right of it; « W. M. C. A. » is as written}

\hand{jo}{Jo Davedson - 80 W. 40" \quad Lo. 5-7233. \\
Mr. Herzberg - City Editor Tribune \\
Paul Beckley - Tribune}

\marginal{incorporate the facts \quad --- up the fore-edge, the pad turned}

\entry{Feb. 13", 47.}{1947-02-13}
\hand{jo}{John Morron -}

\marginal{because in a neighborhood, we can none of us live to ourselves alone
\quad --- at the head of the lower side, to the right of the date}

\hand{jo}{You might like to know what is happening to \\
your neighbors. The people of \#3 are being --- that \\
famous old place that has housed so many \\
distinguished painters + writers --- are being \\
evicted by the general usurper, the N.Y.U. \\
We rejoice that you hold a lease on \#6 + hope \\
the N.Y.U. is not trying to smoke you out. \\
We of \#3 are drawn up for battle against \\
this invasion. It is \struck{pretty} generally known \\
that the Univ. wants to gobble up all the \\
property facing the Sq. then have \struck{their} own lovely \\
Park for its Campus. \struck{They} have all of the \\
E. side + large holdings on the S. If they \\
can oust you in \#6 + prevail on Mrs. \uncertain{Steward} \\
in \#4 to retreat, they have hopes \struck{of} - in time \\
of manoeuvering to get 7-13. - Whereby owning \\
all the best part of Wash. Sq. N. to 5" Ave. \\
\qquad We of \#3 are trying to halt their triumph- \\
ant progress + save Wash. Sq. Park for John Doe \\
+ family + \struck{\#3} the fine studios of \#3 for \\
ourselves. My husband has spent most of his \\
life on the top floor there - 33 1/2 yrs. = 1/3 century. \\
Fred. W. Stokes - now 88 has been there for almost \\
1/2 century. Everyone here had hoped to live + die is this}

\marginal{now \quad --- above « It is generally known »}

\marginal{which is \quad --- above the « = » in « 33 1/2 yrs. = 1/3 century »}

\note{« Fred. W. Stokes » is underlined; « is this » at the end is as written,
and the sentence finishes on page 14; « triumphant » turns the line}

\section{Page 14}

\sheet{96064}{14}
\hand{jo}{old house - often cussed but forever adored. \\
I've always been so fond of your lovely lions* - \\
I'm sure you are lion hearted + won't be shoved \\
about. I have the portrait of one of your lions in \\
my studio, he'll look out for me + mine - \\
\qquad We \struck{meet} of \#3 meet the N.Y.U. aggressors \\
at the O.P.A. next Tues. Feb. 18" \struck{at} \uncertain{11}. \\
We, to show wherefore the N.Y.U. should not evict \\
14 creative artists from a house that for a \\
century has been for artists, many very \\
distinguished. Incl. my husband \struck{\ill{} house has}. \\
The house has produced 9 members of the \\
Institute of Arts + Letters. Can't find one \\
from that N.Y.U. \\
\qquad \qquad Cordially yours.}

\marginal{\uncertain{termite} but gorgeous \quad --- above « old house », at the
head of the page}

\marginal{A.M. \quad --- above « Tues. », with a small « n » below the line}

\marginal{a home \quad --- above « been for artists »}

\hand{jo}{* that you so kindly let me paint some years ago}

\marginal{him \quad --- above « paint », with a caret}

\entry{Feb. 13", 47}{1947-02-13}
\hand{jo}{dear Anita - To thank you for your \struck{kind} heartening letter at Xmas \\
time. \struck{Our case} The case of the N.Y.U. vs. tenants of \#3 - \\
\struck{their aff\ill{}} the combined efforts of \#3 to oppose a certificate \\
for evictions being given to the N.Y.U. has its hearing \\
Tues. A.M. Feb. 18". We all line up with lawyers + \\
petitions signed by the well known painters in the country. \\
The Ass. Press has the story. Feathers flying. \\
\qquad We are asking friends to write to the papers \\
to protest \struck{this deed} the aggression of a \struck{Univ.} \\
supposed to be \struck{interested} in \struck{spreading} Amer. \\
Culture upon a \struck{historic} \struck{building} for a hundred \\
years, has housed people of \struck{long} national distinction. \\
9 members of the Nat. Inst. of Arts + Letters have \\
lived there, incl. E.H. - a present resident. Pictures \\
in museums all over the country have been painted \\
in that house where he has lived for 33 1/2 yrs. \\
F. W. Stokes, who was artist of Staff for Adm. Peary + later \\
accompanied R. Amundsen on \uncertain{arctic} explorations has lived \\
there for \struck{nearly} half a century. Now at 88 he's being \\
evicted - In all 14 creative artists are being sent \\
forth to swell the numbers of displaced persons because \\
a Univ. wishes to expand on \struck{lot of them} + carry their aggression \\
that much farther around Wash. Sq. in an attempt to secure \\
the whole of the Sq. + thus obtain the Park for \struck{its} Campus. \\
We trust the whole country will protest + save the Park for John Doe}

\marginal{wise \quad --- above « kind heartening »}

\marginal{that is \quad --- above « the combined efforts »}

\marginal{by the O.P.A. \quad --- above « given to the N.Y.U. »}

\marginal{a core of \quad --- above « the well known painters », with a caret}

\marginal{the N.Y.U. \quad --- above the struck « Univ. »}

\marginal{an active centre of the arts \quad --- above « a historic building »}

\marginal{they \quad --- above « building », at the right}

\marginal{by him \quad --- above « Pictures »}

\marginal{\uncertain{(4 pages)} \quad --- in the left margin beside « a Univ. wishes
to expand »}

\note{the sentence runs on to page 15. « arctic » is offered against a longer
spelling the letters would also support}

\section{Page 15}

\sheet{96075}{15}
\hand{jo}{+ incidentally, historic \#3 for the artists.}

\hand{jo}{Clayton Bozwell - \\
\qquad Mrs. \uncertain{Roche}.}

\hand{jo}{He would greatly appreciate your signature. I'll explain why later'}

\note{written on the slant, to the right of the two names}

\hand{jo}{Alfred Barr + a number of people on the Mod. Mus. staff \\
also the Whitney, have been signing a petition for \\
E.H. who is about to be evicted from \struck{their} studio \struck{where he} \\
has spent \struck{most of} his life - 1/3 of a century at 3 Wash. Sq. N. \\
He with the rest of the tenants of the house are combatting \\
the aggression of the N.Y.U. who are \struck{\ill{}} trying \\
to get possession of all Wash. Sq. + use the Park for \\
its Campus. They have obtained a lease from the S.S.H. \\
of \#3 + are seeking to evict the tenants, who meet \\
their adversary at 11 am O.P.A. meeting Tues. A.M. Feb. 18. \\
\qquad I, as wife, know what it means to have my husband + \\
his 33 years collection sent adrift in an overcrowded city \\
but this issue is really larger than that. Wash. Sq. has \\
always been perhaps the most distinguished Cultural \\
Centre of the country. These things do not \struck{grow up} readily \\
+ Amer. is not rich in \struck{centres of} Culture. No. 3 is the \\
first opposition the N.Y.U. has encountered in its triumphant \\
sweep. We are told \struck{our} case \struck{legally} is less likely to succeed \\
from the legal angle than from public opinion + that the N.Y.U. \\
will not relish any adverse criticism at this time, because \\
- having already secured 8 000 000 it is now conducting \\
a drive for further contributions for expansion}

\marginal{with the rest of the famous old house \quad --- above « evicted from
their studio »}

\marginal{at 3 W. Sq. \quad --- above « where he », struck}

\marginal{the greater part \quad --- above « most of »}

\marginal{murderously, itself struck \quad --- above the struck word after
« who are »}

\marginal{21 \quad --- above « lease », where the leaf reads « a lease from the
S.S.H. »}

\marginal{with its \struck{well} \quad --- above « Wash. Sq. has »}

\marginal{of art + letters \quad --- above « Cultural Centre »}

\marginal{the famous \#3 \quad --- above « Culture. No. 3 »}

\marginal{our \quad --- above the struck « our »}

\note{the « O » of « O.P.A. » is drawn as a large circle}

\hand{jo}{you will undoubtedly be approached, so the \\
appearance of your name on that petition may give \\
them pause. \struck{Would} Should you think well of letting \\
us have it, Edward H. will be so happy to bring the \\
\uncertain{document} up to you. He has the names of many \\
important painters + yesterday acquired some \\
museum people, so he wouldn't trust that document \\
out of his hands. \qquad If we didn't remember you \\
as one of the sweetest + kindest of \struck{women}, we would \\
venture to \struck{approach} bother you. The loss of that house \\
to the artist would be a calamity, but the \struck{Park} to the \\
people + their baby carriages would be duller \struck{\ill{}} It \\
is very difficult to arouse people until too late. \\
\struck{Please} \struck{add} your \struck{protest} name to our protest - I don't \\
believe a word they say when the N.Y.U. comes to you \\
for a contribution. We know what they'll do with it. \\
\qquad \qquad Earnestly \\
Friedlander. \qquad (740 Park Ave. \\
\qquad \qquad N.Y. 21.)}

\marginal{II. \quad --- at the head of the lower side, at the fore-edge}

\marginal{great ladies \quad --- above « women »}

\marginal{historic \quad --- above « that house »}

\marginal{with splendid studios \quad --- at the head of the following line}

\marginal{still \quad --- above « would be duller »}

\marginal{handsome \quad --- above « for a contribution »}

\hand{jo}{over for I.}

\note{« document » is written twice over itself; the address is drawn round
with an ellipse; the rest of the side is blank. The letter this side ends runs
on from page 16, which is marked I; the two are read 16 then 15, and are given
here in the order the pad is bound}

\section{Page 16}

\sheet{97116}{16}
\hand{jo}{* E.H. has just secured signatures to a petition from \\
Alfred Barr and \struck{of} \struck{several} members of the staff of the Mod. Mus. \\
also of the Whitney. The N.Y.U. has \struck{\ill{}} procured at these \\
a lease of \struck{that} famous old house, where for 100 yrs, they \\
so many distinguished painters + writers have lived. The N.Y.U. is now evicting all present \\
tenants. E.H. has had his studio there for the greater part \\
of his life. An old painter who was artist of Staff for \\
Adm. Peary + is now 88 is also being evicted. In all \\
the 14 creative artists \struck{will} \struck{then} \struck{thrown} \\
will join list of displaced persons turned loose in an \\
over crowded city. \struck{also} it is gradually becoming known that \\
the N.Y.U. \struck{the} hopes to gain possession of all the property \\
surrounding Wash. Sq. + thus obtain the Park for its \\
Campus. It has all the E. side, large holdings on the S. + \\
now starts on that gorgeous row of old white pillared houses \\
on W. S. N. \#3 is the first opposition it has encountered. \\
\qquad The tenants of \#3 + \struck{the} N.Y.U. representative will \\
meet at the O.P.A. next Tues. A.M. We are told \struck{not to} \\
rely \struck{too much from the legal angle}, but our hope comes \\
not from the legal angle but from \uncertain{peoples} opinion - \\
and that the N.Y.U. will not relish adverse criticism \\
because it has a drive on for contributions for \\
expansion. They will undoubtedly approach you, so \\
we feel the presence of your name on that petition \\
will might discouragement in \struck{them}, you will}

\marginal{3 Wash. Sq. N. \quad --- above « famous old house »}

\marginal{there in fine studios \quad --- above « have lived »}

\marginal{nearly 33 1/2 yrs \quad --- above « of his life »}

\marginal{on trips to the N. Pole \quad --- above « Adm. Peary »}

\marginal{the \quad --- above « pillared houses »}

\marginal{on \quad --- above « the 14 creative artists »}

\marginal{be cause \quad --- above « discouragement », with a caret under
« might »}

\marginal{Mrs. \uncertain{Roche} \qquad Feb. 14", \uncertain{47} \quad --- up the
fore-edge, the pad turned}

\hand{jo}{understand where some of that money is used - \\
\struck{to \ill{} they an institutions supposed to at least} \\
\struck{One would expect a University to at least} endorse credible \\
Amer. Art. We believe it ill befits a University at Law \\
to endorse Amer. Culture, not destroy the source of \\
much that is credited at 3 Wash. Sq. And the Sq. so long \\
considered a centre of so much of interest. And the poor Sq. \\
doomed to be swallowed up by the aggressor.}

\marginal{th \quad --- above « at least »}

\marginal{\uncertain{\& endowed} \quad --- above « a University at Law »}

\marginal{to \quad --- above « not destroy », with a caret}

\marginal{highly \quad --- above « is credited »}

\marginal{itself \quad --- at the fore-edge of the « poor Sq. » line}

\hand{jo}{This begins on this page marked I \\
+ continues on page in back " II.}

\note{« marked I » and « II » are underlined. The last two lines are her own
direction for reading the letter: this page first, then the lower side of page
15, which is the physically earlier leaf}

\section{Page 17}

\sheet{97398}{17}
\hand{jo}{LaGuardia \qquad Re Wash. Sq. vs. N.Y.U. \\
\qquad \qquad artists}

\hand{jo}{\qquad As a former trustee of the S.S.H. you must know \\
\struck{the depths} ignominy they can \struck{perpetrate} + also. \\
we all know your heart of gold including humanity \\
as a whole belongs in part to the artists + also \\
baby carriages in Wash. Sq. \\
\qquad Last July, when no one was looking, the S.S.H. \\
sold that famous old 3 Wash. Sq. N. in that row \\
of lovely white pillared red brick houses, the most beautiful \\
block in the city) to the N.Y.U. - ostensibly for \\
administration purposes - for the \struck{18 +} 23. 000 \\
\struck{what} few veterans - who could as readily be \\
dealt with to the E. or the S. of the \struck{Sq.} 14 creative \\
artists are being evicted. \\
\qquad Wash. Sq. has always been perhaps the centre the \\
as the home of its in art + letters in the country. \\
these centres do not spring up with the greed of \\
real estate \struck{\ill{}}. \#3 has a list so long I'll \\
include it at the end. It's been a citadel of the art \\
a little Pantheon *. \struck{the} Present tenants include also \\
people of attainment. By the ruthlessness of the N.Y.U. \\
they are to be evicted from their fine studios to join \\
the list of displaced persons. * Painters require space \\
to paint a picture that will carry across a room \\
or a gallery. It is \struck{necessa} the painter must}

\marginal{what \quad --- above « the depths »}

\marginal{low that \quad --- above « ignominy »}

\marginal{sink to \quad --- above « perpetrate »}

\marginal{a 21 yr. lease of \quad --- above « that famous old 3 Wash. Sq. N. »}

\marginal{few \quad --- above the struck « what »}

\marginal{so the N.Y.U. is in the row \quad --- above « of the Sq. »}

\marginal{\ill{} \quad --- above « are », with a caret under « artists »; a
short word, struck}

\marginal{most distinguished \quad --- above « the centre »}

\marginal{cultural \quad --- above « these centres »}

\marginal{mongers \quad --- above the struck word after « real estate »}

\marginal{high \quad --- above « people of attainment »}

\marginal{(this labelled technical) \quad --- up the fore-edge, the pad turned}

\note{the asterisk beside « a little Pantheon » and a curve beside « Painters
require space » are in red crayon, and « Painters require space » is underlined
in red; this is the first red anywhere in the notebook. The bracket that opens
« (in that lovely row » is not on the leaf, only the one that closes it}

\hand{jo}{constantly step back to see the effect of his work \\
Also he needs a north light \struck{which is} constant - not \\
light one the \struck{right} side of a sitter in the A.M. + light on \\
\struck{left} side in the P.M. --- or also what is painted in afternoon \\
of sunlight will look far otherwise not so illumined. \\
He needs north light. Space + light. These 2 \\
things are rare finds - they are too expensive, \\
no new studio buildings are being constructed \\
they'd be too expensive for a mere creative \\
artist to \ill{}. There are only some 4 or 5 of the \\
fine old studio buildings left in N.Y. - The artists \\
huddle there. They are \uncertain{freedoms} to their preservation. \\
\qquad A fine example is \#3 Wash. Sq. - with 8 \\
good north light studios + other high ceilings + \\
good light fine for voice culture, music recitals \\
(Paderouski played on the 1" floor front \struck{when} as a \\
young man on 1" trip to America) You should know \\
3 Wash. Sq. - the whole place is soaked with tradition \\
of hardihood + accomplishment. It is frequently cussed + \struck{plenty} but forever adored, \\
\qquad * We now learn that \struck{\ill{}} that it has \ill{} \\
the ambition of the N.Y.U. to completely surround \\
Wash. Sq., thereby acquiring the Park for its Campus. \\
It already has all the E. side of the Sq., large holdings \\
on the S. + now begins at the North End end to gobble \\
up that beautiful N. side. \#3 offers the first opposition \\
it has encountered in its triumphant \struck{\ill{}} aggression. \\
* The artists of \#3 maintain}

\marginal{of distance \quad --- above « to see the effect »}

\marginal{with \quad --- above « + other high ceilings »}

\marginal{the \quad --- above « first opposition »}

\note{a red curve runs down the fore-edge beside the Paderouski lines and
under « soaked with tradition ». The line beginning « of hardihood » is written
over a rubbed-out line that is still legible under it and has not been read;
the passage opening « We now learn » is rubbed and written again, and what
stands between « that it has » and « the ambition » will not be read ---
page 19 writes the same sentence out whole. The asterisk before « We now learn »
is in black. The sentence finishes on page 18}

\section{Page 18}

\sheet{96091}{18}
\hand{jo}{Wash. Sq. is for John Doe, wife + baby carriage + \\
\#3 for the artists. \#3 + all that can be saved. \\
\qquad \struck{The} The tenants of \#3 appear at a hearing \\
at the O.P.A. (\struck{the N.Y.U. at their}) Tues. A.M., Feb. 18" - \\
to prove why they should not be evicted. \\
\struck{I tell} you who cares for artists, + song birds! \\
\struck{they'll fight} - + this \struck{could} n't be \\
done in older country that cherishes \\
culture + the traditions of culture. \struck{they} \\
N.Y.U. \struck{fights} \struck{has its own} is in its own field - \\
legality - manipulation of the \struck{legal} chicanerie \\
\qquad (How altogether heartening it would be \\
if the Little Flower of St. Francis would come \\
join us --- \struck{in a cause \ill{}} \\
\qquad The wolf of Gubio was a much more charming \\
creature to deal with than that N.Y.U.}

\marginal{we maintain \quad --- above « Wash. Sq. », at the head of the page}

\marginal{yet \quad --- above « can be saved »}

\marginal{by the N.Y.U. \quad --- above « be evicted »}

\marginal{+ belle \uncertain{lettre}? \quad --- at the fore-edge, below « song
birds! »}

\marginal{this thrown out \quad --- above « this could n't be »}

\marginal{an \quad --- above « in older country », with a caret}

\marginal{the artists \quad --- above « cherishes », at the fore-edge}

\marginal{manages will fight \quad --- above « culture. they », at the
fore-edge}

\marginal{but \quad --- at the left of « N.Y.U. fights »}

\hand{jo}{Incl. in tenants of \#3 are \\
Fred. W. Stokes, artist of Staff for Adm. Peary in his \\
trips to the N. pole. \struck{\ill{}} lived in \#3 for nearly 1/2 century \\
is now 88 - being evicted. Sick in bed? Not all up \\
+ doing battle. \\
E.H., \struck{has} studio here for \struck{his} 33 1/2 yrs. Canvases in all the big \\
\struck{museums from here to the Pacific}. See Who's Who. \\
Mus. of Mod. Art, Whitney, Metrop. -etc. etc. in in Amer. Art}

\marginal{present \quad --- above « in tenants », with a caret}

\marginal{who has been painting for a lifetime \quad --- above « 33 1/2 yrs.
Canvases », written as two interleaved lines}

\hand{jo}{defending cause of Amer. \struck{art} against agression \\
civilization. It's love to think you know that fine \\
\qquad strong canvases - \\
Walter Pach - most distinguished translator of Elie Faure, \\
many books of art criticism - + a painter. One time \\
art prof. of N.Y.U. \\
Jos. Pollet - landscape - Guggenheim fellowship. \\
Mrs. Norman Van Meter - preparing operatic song birds \\
\uncertain{Gunvor Bul Tielman} - murals in Rockefeller Centre. \\
Etc. Etc. All expecting to live + die in the old house - \\
while upholding its traditions}

\marginal{product \quad --- above the struck « art »}

\marginal{\ill{} \quad --- above « against agression », running into the
fore-edge}

\note{« Gunvor Bul Tielman » is the fourth writing of a name pages 3, 5 and 10
give three other ways, and the second letter of the given name is an a or a u
on every one of them. The rest of this side is blank}

\section{Page 19}

\sheet{97044}{19}
\hand{jo}{LaGuardia \qquad Re \struck{evicted} artists of Wash. Sq. \\
\qquad \qquad vs. N.Y.U. \\
As a former trustee of the S.S.H. you must know \\
\struck{the depths of ignominy that institution can sink}. You've \\
\struck{not there.} + We all know your heart of gold \struck{includes} \\
\struck{deepest} \struck{only} humanity as a whole + includes the artists + baby \\
carriages of Wash. Sq. \\
\qquad Last July, when no one was looking, the S.S.H. \\
sold a 21 yr. lease of that famous old 3 Wash. Sq. N. \\
(in that lovely row of red brick with white portals, the \\
most beautiful block in the city) to the N.Y.U., ostensibly \\
for administration purposes - in behalf of the 18 or 23 000 \\
veterans who could be as provided for \struck{elsewhere} \\
\struck{to the} S. or E. of the Square - + not evict 14 creative \\
artists from their studios for the last hundred years have \\
been their birthright. Wash. Sq. has \struck{always} been the \\
most distinguished centre of the arts + letters of the \\
whole country. These cultural centers do not spring up \\
with the greed of real estate mongers - quite to the \\
contrary. \#3 is an \struck{\ill{}} address renowned \\
for the celebrated people who for 100 yrs. have associated \\
with the old house built to fit the most urgent needs \\
of artists. The \struck{list} is so long, I include it later, \\
the particulars of the \struck{people} I explain farther on. \\
It has been an Amer. Pantheon - a citadel of the arts}

\marginal{artists they'd equate \quad --- to the right of « vs. N.Y.U. »}

\marginal{about them + plenty. However \quad --- above « the depths of ignominy »}

\marginal{to what \quad --- at the left of the same line}

\marginal{last summer \quad --- above « not there »}

\marginal{embraces \quad --- above « includes »}

\marginal{which \quad --- above « as a whole », with a caret}

\marginal{marble \quad --- above « white portals »}

\marginal{well \quad --- above « as provided for »}

\marginal{in blocks \quad --- at the left of « S. or E. of the Square »}

\marginal{fine \quad --- above « their studios »}

\marginal{a \quad --- above « birthright », with a caret}

\marginal{for many years \quad --- above the struck « always »}

\marginal{names of these people, known to us \quad --- above « The list is so
long »}

\marginal{needs of these \quad --- above « the particulars of the people »}

\marginal{also \quad --- above « I explain farther on »}

\marginal{centers \quad --- below « a citadel of the arts », with an arrow up
to it}

\hand{jo}{Present tenants continue the tradition of the house \\
with work of \struck{high} attainment. By the ruthlessness of \\
the N.Y.U. they are to be evicted from their irreplace- \\
able workshops + sent forth to swell the list of \\
"displaced persons". Why? \\
\qquad We now learn that it has long been the \\
ambition of the N.Y.U. to completely surround \\
Wash. Sq., thereby acquiring the Park for its \\
Campus. It already has all the E. side of the Sq., \\
large holdings on the S. + now begins at \\
the N.E. end to gobble up that beautiful N. side. \\
\#3 offers the first opposition it has encountered \\
in its triumphant \struck{sweep} tour of \\
aggression. The tenants of \#3 maintain that \\
a birthright is worth fighting for. \\
that Wash. Sq. belongs to John Doe, his wife + the \\
baby carriage + that \#3 belongs to the artists. \\
The tenants of \#3 appear at a hearing at the \\
O.P.A. Tues. A.M. Feb. 18" to combat eviction by \\
the N.Y.U. It was a rousing, thrilling appeal for Amer. Art. \\
\qquad I tell you who cares for artists + song birds + \\
belle lettre. This brutal assault on the life work \\
of serious creative artists. Could'nt happen in an \\
older country that cherishes \struck{its} culture + the \\
traditions of culture, or the well being of little children \\
who need + enjoy the happy park belonging to +}

\marginal{recognized \quad --- above the struck « high »}

\note{the two lines beginning « that Wash. Sq. belongs » are drawn round with
a box and carry a double vertical rule at the left; « irreplaceable » and
« triumphant » turn the line. The sentence finishes on page 20}

\section{Page 20}

\sheet{96251}{20}
\hand{jo}{enjoyed by the whole community. If the park \\
is completely surrounded + enclosed + monopo- \\
lized, John Doe + his will be made to feel \struck{intruders} \\
\struck{+ required to go elsewhere.} \\
that he has ventured into a restricted area, \\
that he is an intruder - \struck{+ required to go elsewhere.} \\
It is only too likely that usurpers will \\
press every advantage. Now is the time to \\
curb this encroachment, now that the \\
design is made clear. \\
\qquad In this fight the N.Y.U. will be on its own field in \\
the manipulation of legal chicaneries. \\
\qquad (How altogether heartening it would be if \\
the Little Flower of St. Francis would come join the \\
good cause. The wolf of Gubio was a \struck{much nicer} \\
very charming beast \struck{compared with \ill{}} than our \\
aggressor, the N.Y.U.}

\marginal{at the O.P.A. Tues. \quad --- above « In this fight the N.Y.U. »}

\marginal{far more \quad --- above « much nicer »}

\marginal{a dear \quad --- above « compared with »}

\note{« monopolized », « irreplaceable » and « restricted » turn the line}

\hand{jo}{List of former tenants incl. \qquad \qquad I \\
\qquad of \quad present \qquad " \qquad \qquad II \\
Explanation of technical needs of artists. \quad III}

\hand{jo}{(2 pages back.)}

\note{« 2 pages back » is in red crayon and stands under the second and third
lines, which a brace gathers; the roman numerals stand in a column at the
fore-edge}

\hand{jo}{I. Former tenants include - painters + sculptors \\
Wm. M. Glackens* \quad Mr. + Mrs. \\
Guy Pène du Bois * \\
Ernest Lawson * \\
Rockwell Kent. \\
\uncertain{Conceta} Scaravaglioni \\
Brenda Putnam * \\
Marjorie Ryerson \\
\uncertain{Aileen Dresser} \\
Mary Foot \\
Kimon Nicolades, \\
Bobby Edwards \\
Thomas Weber Furlong \\
Stephen Haweis \\
Grace Randolf. \\
Mary Tillinghast \\
\uncertain{Wallis Fitts} \\
\uncertain{Wailes} Goedbeck. \\
Benj. Porter \\
the \uncertain{Monines} * \\
Anne Dunbar.}

\hand{jo}{Alex Brook * \\
Coulton Waugh \\
Ricco LeBrun \\
Waldo Pierce \\
Writers + others incl. \\
John Dos Passos \\
E. E. Cummings \\
Schofield Thayer * \\
Edmund Wilson * \\
Eleanor Wylie \\
Frank Harris \\
F. Hopkinson Smith \\
Mrs. Jack Kelly. \\
Albert White Vorse \\
Lillian + Dorothy Gish \\
" Langhorn etc. etc. - \\
\qquad so many more}

\note{the list runs in two columns, the second beginning at « Alex Brook » and
standing to the right of the first; the two are given here one after the other.
Red crayon marks the left of eleven of the names in the first column with a
dash, sets a red asterisk beside « Wm. M. Glackens », « Guy Pène du Bois »,
« Ernest Lawson », « Brenda Putnam », « the Monines », « Alex Brook »,
« Schofield Thayer » and « Edmund Wilson », and sets a red cross beside
« Writers + others incl. » at the head of the second column. Faint pencil
numerals
stand in the left margin --- « 19 » beside Mary Tillinghast, « 2 » beside
Wailes Goedbeck, « 20 » beside Benj. Porter and « 23 » beside the Monines ---
and their sense is not established. « Conceta Scaravaglioni » is a third
spelling against pages 2, 4 and 9, and the strokes here read as readily
« Coreata »}

\hand{jo}{So many others known to \\
Amer. Art}

\hand{jo}{\qquad Incl. present tenants \quad 9 members of Inst. of Arts + Letters \\
Fred. W. Stokes - 88 - lived in Wash. Sq. N. 45 yrs. - Artist of Staff of \\
\qquad Adm. Peary on trips to N. Pole. \\
E. H. * - \struck{has} lived here 33 1/2 yrs. the \struck{\ill{}} of his life. Canvases \\
in all the big museums \struck{from coast to coast} in the Country \\
\qquad a staunch defender of the art of America \\
Walter Pach - painter - + critic - Beautiful translation of \\
\qquad Elie Faure. Hist of Art in 3 vols. \\
Jos. Pollet - painter - \uncertain{Guggenheimes} fellowship - teaching \\
\qquad art in N.Y.U. \\
Mrs. Norman Van Meter - preparing operatic song birds under her \\
\qquad fine high ceilings}

\note{the asterisk beside « E. H. » is in black}

\section{Page 21}

\sheet{97153}{21}
\hand{jo}{\uncertain{Louro More} - influencing " modern jewelry' - \\
\qquad etc. etc. \qquad All the wives good painters. \\
7 veterans, incl. 2 colonels.}

\marginal{an \quad --- above « influencing »}

\hand{jo}{Some 10 yrs. ago the N.Y.U. threw John Sloan out of \\
his fine studio at the Judson. A brutal thing to \\
do. It dried up the fine stream of canvases depicting \\
Wash. Sq., \struck{the park} + the life of this region - a fine \\
historic statement -. \\
\qquad A trustee of the Mus. of Modern Art fears this is what \\
will happen to E. Hopper. How I hope you know the \\
work of my husband? It just occurs to me to \struck{bring} \\
one of his more recent books. Some time read his Credo \\
He's like that - good, honest to the core - as lots a \\
fine sensitive artist - too sensitive - I fight to \\
get this book - he says you don't want his books, at least \\
that we were trying to influence you - What \struck{else} \\
ever else? We'd be so happy to feel you could \\
be with us - advise us wisely - because we all \\
lean on you! \qquad Earnestly - Jo N. H. \\
\qquad \qquad \qquad Mrs. Edw}

\marginal{its \quad --- above « the park », with a caret}

\marginal{American \quad --- below « historic statement »}

\marginal{send you \quad --- above « to bring »}

\marginal{own \quad --- above « his Credo »}

\hand{unidentified}{Fiora}

\note{« Fiora » stands alone in the blank lower half of the side, in a much
darker and blunter medium than the pencil of everything round it, and the hand
cannot be placed. The signature above it trails off after « Mrs. Edw »}

\entry{Feb. 20" Thurs.}{1947-02-20}
\hand{jo}{\uncertain{\&} Mayor LaGuardia -- \\
\qquad As a former trustee of the Sailors Snug Harbor, \\
you know what they're like. However you weren't there \\
last summer. We all know your heart of gold \\
embraces all decent humanity, which includes \\
the artists. + baby carriages. \\
\qquad Last July, when no one was looking, the Sailors \\
Snug sold a 21 yr. lease of that famous old 3 Wash. \\
in that lovely row of red brick with white \struck{marble} \\
columns, the most beautiful block in the city) to N.Y.U., \\
ostensibly for administration purposes - in behalf of \\
swarms of veterans - who could be as well, far better, \\
provided for \struck{the} in blocks S. + E. of the Square - \\
+ not evict 14 creative artists from their fine studios \\
that for the last century have been their birthright. \\
- studios rare, no new ones being built. Wash. Sq. \\
for so long most distinguished centre of arts + letters \\
in whole country. (I need'nt try to sell you Wash. Sq.) \\
These cultural centers do not spring up with the greed \\
of real estate mongers - quite the contrary. \\
\#3 is an address renowned - even in Europe - for the \\
celebrated people associated with the old house built to \\
fit their special needs - of space + north light - the list so \\
long I enclose it later. The house has been an Amer. Pan- \\
theon - historic like Concord + the transcendentalists}

\marginal{God love them \quad --- above « who could be as well », with a caret
under « veterans »}

\note{the opening mark before « Mayor » is a single character that will not be
read. The sentence finishes on page 22}

\section{Page 22}

\sheet{97309}{22}
\hand{jo}{Present tenants continue the tradition with \\
works of recognized attainment. By the ruthless- \\
ness of N.Y.U. they are to be evicted from their irre- \\
placeable workshops (+ dwellings) + sent forth to swell \\
the list of displaced persons. Why? \\
\qquad We now learn that it has long been the ambition \\
of the N.Y.U. to completely surround Wash. Sq., thereby \\
securing the Park for its Campus. It already has \\
all the E. side, considerable holdings on the S., + \\
now begins + \uncertain{gains} to gobble up the N. side. \\
\#3 offers the first opposition it has encountered in its \\
triumphant tour of aggression. The artists \\
feel a birthright is worth a battle. \\
\qquad There was a fine rousing meeting at the O.P.A. \\
Mrs. John Sloan + Walter Pach were magnificent - \\
such stark sincerity + nobility - + our polar \\
explorer aged 88! \\
\qquad I tell you, who needs no telling, this brutal \\
attack on the life work of serious creative artists \\
could'nt happen in an older country that cherishes \\
its own culture, or the welfare of little children \\
who need + enjoy the happy park enjoyed by the whole \\
community. If the park is completely surround- \\
ed - enclosed + monopolized, John Doe, wife \\
+ baby carriage will be made to feel he has ventured}

\marginal{beautiful \quad --- above « the N. side », with a caret}

\hand{jo}{into a restricted area + is an intruder + \\
required to go elsewhere. \\
\qquad In this fight \struck{the} N.Y.U. will be on its own field \\
manipulation of legal chicaneries. \\
\qquad How altogether heartening it would be if the \\
Little Flower of St. Francis would join us. \\
The wolf of Gubio was a far more charming \\
beastie to deal with than that aggressive N.Y.U. \\
historic}

\marginal{the \quad --- at the fore-edge of « its own field »}

\hand{jo}{List of former tenants, back 2 pages. \\
\qquad " \quad " \quad present \qquad " \quad "}

\note{« historic » is struck; the rest of this side is blank}

\section{Page 23}

\sheet{96094}{23}
\entry{Tues. P.M. Feb. 25", 47.}{1947-02-25}
\hand{jo}{Invited to party \\
\qquad at \uncertain{E. Dunscombe Marie's.}}

\hand{jo}{Guy du Bois * \quad Mr. + Mrs. Yvonne, \uncertain{Hopkins} \\
\qquad \qquad Billy - Jane - \\
Frank + Margaret Kelly \\
Sarah Newmeyer 37 W. 8 - \\
Harvey Wiley Corbett. Mr. + Mrs. \quad 3rd + 4" ave. \quad Sp. 7-7680 \\
Reggie Marsh * + Felicea \\
James Chapin - Mr. + Mrs. \\
Pierce W. Walters \quad " \quad " \quad + Miss \\
Morris Ernst \\
Eric Gugler * 101 Park Ave. \\
Deems Taylor * \quad Mr. Mrs. \quad 2 E. 66" \\
Douglas Moore * \quad Mr. \quad " \\
Austin Strong * \quad " \quad 707 Mad. Ave. \quad Re. 7-0434 \\
\uncertain{Stark} Young * \quad " \quad " \quad 320 E. 57" \\
Gifford Beal * \quad " \quad 27 W. 67" \\
Henry Sudd Canby * ed. \quad " \quad 25 W. 45" \\
Mrs. Force - phone. \\
John Sloan * + Helen (\uncertain{Alogan}) \\
Everett Shinn \quad " \\
St. Gaudens - phone. no cando. \\
Whitelaw Reid - letter. Mrs. right - taken \\
Mr. Mrs. Arthur Paul Waltheim \\
Mr. + Mrs. Maurice Whitebook \\
Emma Bugbee}

\hand{jo}{Leon Kroll 15 W. 67"}

\marginal{No cando \quad --- above « Leon Kroll », at the fore-edge}

\marginal{teaching \quad --- to the right of « James Chapin », below a struck
« maybe » that stands beside « Felicea »}

\marginal{Hauptman \quad --- to the right of « Waltheim »}

\marginal{partners. \quad --- to the right of « Whitebook », below « Hauptman »,
the two braced to the two couples}

\note{a column of marks stands at the left of the names --- a tick against
Guy du Bois, Frank + Margaret Kelly, Sarah Newmeyer, James Chapin, Morris Ernst,
Eric Gugler and John Sloan; a cross against Harvey Wiley Corbett, Leon Kroll,
Austin Strong, Henry Sudd Canby, St. Gaudens and Whitelaw Reid; a query against
Gifford Beal and Mrs. Force. « letter » or « letters » is written above nine of
the lines and « " » repeats it down the column. Red underlines run under
« du Bois », « Newmeyer », « Corbett », « Ernst » and « Force », and red
asterisks stand beside several of the names}

\hand{jo}{Mrs. Glackens. \\
Mark Van Doren * 393 Bleaker \\
\qquad Ch. 3 - 7441 \qquad out of town \quad don't care. \\
Mrs. deVilbis 28 E. 10" \\
\qquad Sec. Wash. Sq. Neighbors \\
\qquad Gr. 3 - 4173 \\
\struck{Carl Van Doren * 41 Cents. Pk. W.} \\
\qquad Tr. 7 - 8500 \\
Mr. Mrs. Chas. \uncertain{Chap.} Keck 40 W. 10" \\
\qquad St. 9 - 1236 \\
\struck{Judge Jules Isaacs} - 21 E. 10" \\
\qquad Al. 4 - 5713. \\
Morris Ernst \qquad W. 11" \\
Mrs. Wm. Glackens 10 W. 9" \\
\qquad St. 9 - 0873. \\
Shotwell of Columbia. \\
Moe - Guggenheim}

\marginal{Pollet \quad --- above « Guggenheim »}

\marginal{out of town \quad --- above « Guggenheim », to the left of « Pollet »}

\note{a cross stands against Mark Van Doren, Shotwell and Moe, and a tick
against Judge Jules Isaacs; the same red and black marking runs down this side
as down the upper one. The list is the whole of the lower side and the sitting
ends here}

\keywords{Washington Square, New York University, Sailors Snug Harbor,
3 Washington Square North, John Sloan, evictions, Office of Price
Administration, Herald Tribune, Greenwich Village associations, studio houses,
Frederick W. Stokes, drafts of letters, Fiorello LaGuardia, north light,
tenant lists, Walter Pach, petitions}
\end{document}
